\documentclass[a4paper,12pt]{article}
\usepackage[T1]{fontenc}
\usepackage[german]{babel}

\renewcommand{\leq}{\mathrel{\vcenter{\halign{\hfil$##$\hfil\cr
   =\cr\noalign{\kern-8pt}<\cr}}}}

\renewcommand{\geq}{\mathrel{\vcenter{\halign{\hfil$##$\hfil\cr
   >\cr\noalign{\kern-8pt}=\cr}}}}


\newcommand{\asterisktriangle}{\centerline{${\displaystyle
   * \hskip 36pt * \atop \displaystyle *}$}}

\begin{document}

\thispagestyle{empty}
\begin{center}
\Large\bfseries
Grundlagen f\"{u}r eine allgemeine Theorie der Functionen
einer ver\"{a}nderlichen complexen Gr\"{o}sse.\\[12 pt]
Bernhard Riemann\\[12 pt]
[Inauguraldissertation, G\"{o}ttingen, 1851; zweiter
unver\"{a}nderter Abdruck, G\"{o}ttingen 1867.]\\[24 pt]
\large\mdseries
Transcribed by D. R. Wilkins\\[12 pt]
Preliminary Version: December 1998\\
Corrected: April 2000
\end{center}

\newpage
\setcounter{page}{1}

\title{Grundlagen f\"{u}r eine allgemeine Theorie der Functionen
einer ver\"{a}nderlichen complexen Gr\"{o}sse.}
\author{Bernhard Riemann}
\date{[Inauguraldissertation, G\"{o}ttingen, 1851; zweiter
unver\"{a}nderter Abdruck, G\"{o}ttingen 1867.]}

\maketitle

\centerline{1.}

\nobreak\medskip

Denkt man sich unter $z$ eine ver\"{a}nderliche Gr\"{o}sse,
welche nach und nach alle m\"{o}glichen reellen Werthe annehmen
kann, so wird, wenn jedem ihrer Werthe ein einziger Werth der
unbestimmten Gr\"{o}sse $w$ entspricht, $w$ eine Function von $z$
genannt, und wenn, w\"{a}hrend $z$ alle zwischen zwei festen
Werthen gelegenen Werthe stetig durchl\"{a}uft, $w$ ebenfalls
stetig sich \"{a}ndert, so heisst diese Function innerhalb dieses
Intervalls stetig oder continuirlich.

Diese Definition setzt offenbar zwischen den einzelnen Werthen
der Function durchaus kein Gesetz fest, indem, wenn \"{u}ber
diese Function f\"{u}r ein bestimmtes Intervall verf\"{u}gt ist,
die Art ihrer Fortsetzung ausserhalb desselben ganz der
Willk\"{u}r \"{u}berlassen bleibt.

Die Abh\"{a}ngigkeit der Gr\"{o}sse $w$ von $z$ kann durch ein
mathematisches Gesetz gegeben sein, so dass durch bestimmte
Gr\"{o}ssenoperationen zu jedem Werthe von $z$ das ihm
entsprechende $w$ gefunden wird.  Die F\"{a}higkeit, f\"{u}r alle
innerhalb eines gegebenen Intervalls liegenden Werthe von $z$
durch dasselbe Abh\"{a}ngigkeitsgesetz bestimmt zu werden,
schrieb man fr\"{u}her nur einer gewissen Gattung von Functionen
zu (functiones continuae nach Euler's Sprachgebrauch); neuere
Untersuchungen haben indess gezeigt, dass es analytische
Ausdr\"{u}cke giebt, durch welche eine jede stetige Function
f\"{u}r ein gegebenes Intervall dargestellt werden kann.  Es ist
daher einerlei, ob man die Abh\"{a}ngigkeit der Gr\"{o}sse $w$
von der Gr\"{o}sse $z$ als eine willk\"{u}rlich gegebene oder als
eine durch bestimmte Gr\"{o}ssenoperationen bedingte definirt.
Beide Begriffe sind in Folge der erw\"{a}hnten Theoreme congruent.

Anders verh\"{a}lt es sich aber, wenn die Ver\"{a}nderlichkeit
der Gr\"{o}sse $z$ nicht auf reelle Werthe beschr\"{a}nkt wird,
sondern auch complexe von der Form $x + yi$ (wo $i = \sqrt{-1}$)
zugelassen werden.

Es seien $x + yi$ und $x + yi + dx + dy \, i$ zwei unendlich
wenig verschiedene Werthe der Gr\"{o}sse $z$, welchen die Werthe
$u + vi$ und $u + vi + du + dv \, i$ der Gr\"{o}sse $w$
entsprechen.  Alsdann wird, wenn die Abh\"{a}ngigkeit der
Gr\"{o}sse $w$ von $z$ eine willk\"{u}rlich angenommenene ist,
das Verh\"{a}ltniss
$\displaystyle\frac{du + dv \, i}{dx + dy \, i}$
sich mit den Werthen von $dx$ und $dy$, allgemein zu reden,
\"{a}ndern, indem, wenn man
$dx + dy \, i = \varepsilon e^{\varphi i}$
setzt,
\begin{eqnarray*}
\frac{du + dv \, i}{dx + dy \, i}
   &=&   {\textstyle\frac{1}{2}}
         \left(
            \frac{\partial u}{\partial x}
          + \frac{\partial v}{\partial y}
         \right)
       + {\textstyle\frac{1}{2}}
         \left(
            \frac{\partial v}{\partial x}
          - \frac{\partial u}{\partial y}
         \right) i \\
   & & + {\textstyle\frac{1}{2}}
         \left[
            \frac{\partial u}{\partial x}
          - \frac{\partial v}{\partial y}
          + \left(
               \frac{\partial v}{\partial x}
             + \frac{\partial u}{\partial y}
            \right) i
         \right]
         \frac{dx - dy \, i}{dx + dy \, i} \\
   &=&   {\textstyle\frac{1}{2}}
         \left(
            \frac{\partial u}{\partial x}
          + \frac{\partial v}{\partial y}
         \right)
       + {\textstyle\frac{1}{2}}
         \left(
            \frac{\partial v}{\partial x}
          - \frac{\partial u}{\partial y}
         \right) i \\
   & & + {\textstyle\frac{1}{2}}
         \left[
            \frac{\partial u}{\partial x}
          - \frac{\partial v}{\partial y}
          + \left(
               \frac{\partial v}{\partial x}
             + \frac{\partial u}{\partial y}
            \right) i
         \right]
         e^{-2 \varphi i}
\end{eqnarray*}
wird.  Auf welche Art aber auch $w$ als Function von $z$ durch
Verbindung der einfachen Gr\"{o}ssenoperationen bestimmt werden
m\"{o}ge, immer wird der Werth des Differentialquotienten
$\displaystyle \frac{dw}{dz}$
von dem besondern Werthe des Differentials $dz$ unabh\"{a}ngig
sein\footnote{Diese Behauptung ist offenbar in allen F\"{a}llen
gerechtfertigt, wo sich aus dem Ausdrucke von $w$ durch $z$
mittelst der Regeln der Differentiation ein Ausdruck von
$\displaystyle \frac{dw}{dz}$
durch $z$ finden l\"{a}sst; ihre streng allgemeine G\"{u}ltigkeit
bleibt f\"{u}r jetzt dahin gestellt.}.  Offenbar kann also auf
diesem Wege nicht jede beliebige Abh\"{a}ngigkeit der complexen
Gr\"{o}sse $w$ von der complexen Gr\"{o}sse $z$ ausgedr\"{u}ckt
werden.

Das eben hervorgehobene Merkmal aller irgendwie durch
Gr\"{o}ssenoperationen bestimmbaren Functionen werden wir f\"{u}r
die folgende Untersuchung, wo eine solche Function unabh\"{a}ngig
von ihrem Ausdrucke betrachtet werden soll, zu Grunde legen,
indem wir, ohne jetzt dessen Allgemeing\"{u}ltigkeit und
Zul\"{a}nglichkeit f\"{u}r den Begriff einer durch
Gr\"{o}ssenoperationen ausdr\"{u}ckbaren Abh\"{a}ngigkeit zu
beweisen, von folgender Definition ausgehen:

Eine ver\"{a}nderliche complexe Gr\"{o}sse $w$ heisst eine
Function einer andern ver\"{a}nderlichen complexen Gr\"{o}sse
$z$, wenn sie mit ihr sich so \"{a}ndert, dass der Werth des
Differentialquotienten
$\displaystyle \frac{dw}{dz}$
unabh\"{a}ngig von dem Werthe des Differentials $dz$ ist.

\medbreak

\centerline{2.}

\nobreak\medskip

Sowohl die Gr\"{o}sse $z$, als die Gr\"{o}sse $w$ werden als
ver\"{a}nderliche Gr\"{o}ssen betrachtet, die jeden complexen
Werth annehmen k\"{o}nnen.  Die Auffassung einer solchen
Ver\"{a}nderlichkeit, welche sich auf ein zusammenh\"{a}ngendes
Gebiet von zwei Dimensionen erstreckt, wird wesentlich
erleichtert durch eine Ankn\"{u}pfung an r\"{a}umliche
Anschauungen.

Man denke sich jeden Werth $x + yi$ der Gr\"{o}sse $z$
repr\"{a}sentirt durch einen Punkt~$O$ der Ebene~$A$, dessen
rechwinklige Coordinaten $x$,~$y$, jeden Werth $u + vi$ der
Gr\"{o}sse $w$ durch einen Punkt~$Q$ der Ebene~$B$, dessen
rechtwinklige Coordinaten $u$, $v$ sind.  Eine jede
Abh\"{a}ngigkeit der Gr\"{o}sse $w$ von $z$ wird sich dann
darstellen als eine Abh\"{a}ngigkeit der Lage des Punktes~$Q$ von
der des Punktes~$O$.  Entspricht jedem Werthe von $z$ ein
bestimmter mit $z$ stetig sich \"{a}ndernder Werth von $w$, mit
andern Worten, sind $u$ und $v$ stetige Functionen von $x$,~$y$,
so wird jedem Punkte der Ebene~$A$ ein Punkt der Ebene~$B$, jede
Linie, allgemein zu reden, eine Linie, jedem
zusammenh\"{a}ngenden Fl\"{a}chenst\"{u}cke ein
zusammenh\"{a}ngendes Fl\"{a}chenst\"{u}ck entsprechen.  Man wird
sich also diese Abh\"{a}ngigkeit der Gr\"{o}sse $w$ von $z$
vorstellen k\"{o}nnen als eine Abbildung der Ebene~$A$ auf der
Ebene~$B$.

\medbreak

\centerline{3.}

\nobreak\medskip

Es soll nun untersucht werden, welche Eigenschaft diese Abbildung
erh\"{a}lt, wenn $w$ eine Function der complexen Gr\"{o}sse $z$,
d.~h.\ wenn
$\displaystyle \frac{dw}{dz}$
von $dz$ unabh\"{a}ngig ist.

Wir bezeichnen durch $o$ einen unbestimmten Punkt der Ebene~$A$
in der N\"{a}he von $O$, sein Bild in der Ebene $B$ durch $q$,
ferner durch
$x + yi + dx + dy \, i$ und
$u + vi + du + dv \, i$ die Werthe der Gr\"{o}ssen $z$ und $w$ in
diesen Punkten.  Es k\"{o}nnen dann $dx$, $dy$ und $du$, $dv$ als
rechtwinklige Coordinaten der Punkte $o$ und $q$ in Bezug auf die
Punkte $O$ und $Q$ als Anfangspunkte angesehen werden, und wenn
man
$dx + dy \, i = \varepsilon e^{\varphi i}$
und
$du + dv \, i = \eta e^{\psi i}$
setzt, so werden die Gr\"{o}ssen
$\varepsilon$, $\phi$, $\eta$, $\psi$ Polarcoordinaten dieser
Punkte f\"{u}r dieselben Anfangspunkte sein.  Sind nun $o'$ und
$o''$ irgend zwei bestimmte Lagen des Punktes $o$ in unendlicher
N\"{a}he von $O$, und dr\"{u}ckt man die von ihnen abh\"{a}ngigen
Bedeutungen der \"{u}brigen Zeichen durch entsprechene Indices
aus, so giebt die Voraussetzung
\[    \frac{du'  + dv'  \, i}{dx'  + dy'  \, i}
    = \frac{du'' + dv'' \, i}{dx'' + dy'' \, i} \]
und folglich
\[    \frac{du'  + dv'  \, i}{du'' + dv'' \, i}
    = \frac{\eta'}{\eta''} e^{(\psi' - \psi'') i}
    = \frac{dx'  + dy'  \, i}{dx'' + dy'' \, i}
    = \frac{\varepsilon'}{\varepsilon''} e^{(\varphi' - \varphi'') i},\]
woraus
$\displaystyle
      \frac{\eta'}{\eta''}
    = \frac{\varepsilon'}{\varepsilon''}$
und
$\psi' - \psi'' = \varphi' - \varphi''$,
d.~h.\ in den Dreiecken $o' O o''$ und $q' Q q''$ sind die Winkel
$o' O o''$ und $q' Q q''$ gleich und die sie einschliessenden
Seiten einander proportional.

Es findet also zwischen zwei einander entsprechenden unendlich
kleinen Dreiecken und folglich allgemein zwischen den kleinsten
Theilen der Ebene~$A$ und ihres Bildes auf der Ebene~$B$
Aehnlichkeit Statt.  Eine Ausnahme von diesem Satze tritt nur in
den besonderen F\"{a}llen ein, wenn die einander entsprechenden
Aenderungen der Gr\"{o}ssen $z$ und $w$ nicht in einem endlichen
Verh\"{a}ltnisse zu einander stehen, was bei Herleitung desselben
stillschweigend vorausgesetzt ist\footnote{Ueber diesen
Gegenstand sehe man:

\glqq Allgemeine Aufl\"{o}sung der Aufgabe: Die Theile einer
gegebenen Fl\"{a}che so abzubilden, dass die Abbildung dem
Abgebildeten in den kleinsten Theilen \"{a}hnlich wird\grqq,
von \emph{C.~F.~Gauss}.  (Als Beantwortung der von der
k\"{o}niglichen Societ\"{a}t der Wissenschaften in Copenhagen
f\"{u}r 1822 aufgegebenen Preisfrage, abgedruckt in: \glqq
Astronomische Abhandlungen, herausgegeben von Schumacher.
Drittes Heft. Altona. 1825.\grqq)  (Gauss Werke Bd.~IV, p.~189.)}.

\medbreak

\centerline{4.}

\nobreak\medskip

Bringt man den Differentialquotienten
$\displaystyle \frac{du + dv \, i}{dx + dy \, i}$
in die Form
\[ \frac{\displaystyle
         \left(
            \frac{\partial u}{\partial x}
          + \frac{\partial v}{\partial x} i
         \right) \, dx
       + \left(
            \frac{\partial v}{\partial y}
          - \frac{\partial u}{\partial y} i
         \right) \, dy \, i}{dx + dy \, i},\]
so erhellt, dass er und zwar nur dann f\"{u}r je zwei Werthe von
$dx$ und $dy$ denselben Werth haben wird, wenn
\[       \frac{\partial u}{\partial x}
   =     \frac{\partial v}{\partial y}
      \quad \mbox{und} \quad
         \frac{\partial v}{\partial x}
   =   - \frac{\partial u}{\partial y} \]
ist.  Diese Bedingungen sind also hinreichend und nothwendig,
damit $w = u + vi$ eine Function von $z = x + yi$ sei.  F\"{u}r
die einzelnen Glieder dieser Function fliessen aus ihnen die
folgenden:
\[    \frac{\partial^2 u}{\partial x^2}
    + \frac{\partial^2 u}{\partial y^2}
   = 0,\quad
      \frac{\partial^2 v}{\partial x^2}
    + \frac{\partial^2 v}{\partial y^2}
   = 0,\]
welche f\"{u}r die Untersuchung der Eigenschaften, die Einem
Gliede einer solchen Function einzeln betrachtet zukommen, die
Grundlage bilden.  Wir werden den Beweis f\"{u}r die wichtigsten
dieser Eigenschaften einer eingehenderen Betrachtung der
vollst\"{a}ndigen Function voraufgehen lassen, zuvor aber noch
einige Punkte, welche allgemeineren Gebieten angeh\"{o}ren,
er\"{o}rten und festlegen, um uns den Boden f\"{u}r jene
Untersuchungen zu ebenen.

\nobreak\medskip

\asterisktriangle

\medbreak

\centerline{5.}

\nobreak\medskip

F\"{u}r die folgenden Betrachtungen beschr\"{a}nken wir die
Ver\"{a}nderlichkeit der Gr\"{o}ssen $x$,~$y$ auf ein endliches
Gebiet, indem wir als Ort des Punktes $O$ nicht mehr die Ebene
$A$ selbst, sondern eine \"{u}ber dieselbe ausgebreitete
Fl\"{a}che $T$ betrachten.  Wir w\"{a}hlen diese Einkleidung, bei
der es unanst\"{o}ssig sein wird, von auf einander liegenden
Fl\"{a}chen zu reden, um die M\"{o}glichkeit offen zu lassen,
dass der Ort des Punktes~$O$ \"{u}ber denselben Theil der
Ebene sich mehrfach erstrecke, setzen jedoch f\"{u}r einen solchen
Fall voraus, dass die auf einander liegenden Fl\"{a}chentheile
nicht l\"{a}ngs einer Linie zusammenh\"{a}ngen, so dass eine
Umfaltung der Fl\"{a}che, oder eine Spaltung in auf einander
liegende Theile nicht vorkommt.

Die Anzahl der in jedem Theile der Ebene auf einander liegenden
Fl\"{a}chentheile ist alsdann vollkommen bestimmt, wenn die
Begrenzung der Lage und dem Sinne nach (d.~h.\ ihre innere und
\"{a}ussere Seite) gegeben ist; ihr Verlauf kann sich jedoch noch
verschieden gestalten.

In der That, ziehen wir durch den von der Fl\"{a}che bedeckten
Theil der Ebene eine beliebige Linie~$l$, so \"{a}ndert sich die
Anzahl der \"{u}ber einander liegenden Fl\"{a}chentheile nur beim
Ueberschreiten der Begrenzung, und zwar beim Uebertritt von
Aussen nach Innen um $+1$, im entgegengesetzten Falle um $-1$,
und ist also \"{u}berall bestimmt.  L\"{a}ngs der Ufers dieser
Linie setzt sich nun jeder angrenzende Fl\"{a}chentheil auf ganz
bestimmte Art fort, so lange die Linie die Begrenzung nicht
trifft, da eine Unbestimmtheit jedenfalls nur in einem einzelnen
Punkte und also entweder in einem Punkte der Linie selbst oder in
einer endlichen Entfernung von derselben Statt hat; wir
k\"{o}nnen daher, wenn wir unsere Betrachtung auf einen im Innern
der Fl\"{a}che verlaufenden Theil der Linie~$l$ und zu beiden
Seiten auf einen hinreichend kleinen Fl\"{a}chenstreifen
beschr\"{a}nken, von \emph{bestimmten} angrenzenden
Fl\"{a}chentheilen reden, deren Anzahl auf jeder Seite gleich
ist, und die wir, indem wir der Linie eine bestimmte Richtung
beilegen, auf der Linken mit $a_1, a_2,\ldots\, a_n$, auf der
Rechten mit $a'_1, a'_2,\ldots\, a'_n$ bezeichnen.  Jeder
Fl\"{a}chentheil~$a$ wird sich dann in einen der
Fl\"{a}chentheile $a'$ fortzetzen; dieser wird zwar im
Allgemeinen f\"{u}r den ganzen Lauf der Linie~$l$ derselbe sein,
kann sich jedoch f\"{u}r besondere Lagen von $l$ in einem ihrer
Punkte \"{a}ndern.  Nehmen wir an, dass oberhalb eines solchen
Punktes $\sigma$ (d.~h.\ l\"{a}ngs des vorhergehenden Theils von
$l$) mit den Fl\"{a}chentheilen $a'_1, a'_2,\ldots\, a'_n$ der
Reihe nach die Fl\"{a}chentheile $a_1, a_2,\ldots\, a_n$
verbunden seien, unterhalb desselben aber die Fl\"{a}chentheile
$a_{\alpha_1}, a_{\alpha_2},\ldots\, a_{\alpha_n}$,
wo $\alpha_1, \alpha_2,\ldots\, \alpha_n$ nur in der Anordnung
von $1,2,\ldots,n$ verschieden sind, so wird ein oberhalb
$\sigma$ von $a_1$ in $a'_1$ eintretender Punkt, wenn er
unterhalb $\sigma$ auf die linke Seite zur\"{u}cktritt, in den
Fl\"{a}chentheil $a_{\alpha_1}$ gelangen, und wenn er den
Punkt~$\sigma$ von der Linken zur Rechten umkreiset, wird der
Index des Fl\"{a}chentheils, in welchem er sich befindet, der
Reihe nach die Zahlen
\[ 1, \alpha_1, \alpha_{\alpha_1},\ldots\, \mu, \alpha_\mu,\ldots \]
durchlaufen.  In dieser Reihe sind, so lange das Glied~$1$ nicht
wiederkehrt, nothwendig alle Glieder von einander verschieden,
weil einem beliebigen mittlern Gliede $a_\mu$ nothwendig $\mu$
und nach einander alle fr\"{u}heren Glieder bis $1$ in
unmittelbarer Folge vorhergehen; wenn aber nach einer Anzahl von
Gliedern, die offenbar kleiner als $n$ sein muss und $= m$ sei,
das Glied~$1$ wiederkehrt, so m\"{u}ssen die \"{u}brigen Glieder
in derselben Ordnung folgen.  Der um $\sigma$ sich bewegende
Punkt kommt alsdann nach je $m$ Uml\"{a}ufen in denselben
Fl\"{a}chentheil zur\"{u}ck und ist auf $m$ der auf einander
liegenden Fl\"{a}chentheile eingeschr\"{a}nkt, welche sich
\"{u}ber $\sigma$ zu einem einzigen Punkte vereinigen.  Wir
nennen diesen Punkt einen Windungspunkt $(m-1) \,$\-ter Ordnung der
Fl\"{a}che~$T$.  Durch Anwendung desselben Verfahrens auf die
\"{u}brigen $n - m$ Fl\"{a}chentheile werden diese, wenn sie
nicht gesondert verlaufen, in Systeme von $m_1, m_2,\ldots$
Fl\"{a}chentheilen zerfallen, in welchem Falle auch noch
Windungspunkte $(m_1 - 1) \,$\-ter, $(m_2 - 1) \,$\-ter$\ldots$ Ordnung in
dem Punkte~$\sigma$ liegen.

Wenn die Lage und der Sinn der Begrenzung von $T$ und die Lage
ihrer Windungspunkte gegeben ist, so ist $T$ entweder vollkommen
bestimmt oder doch auf eine endliche Anzahl verschiedener
Gestalten beschr\"{a}nkt; Letzteres, in so fern sich diese
Bestimmungsst\"{u}cke auf verschiedene der auf einander liegenden
Fl\"{a}chentheile beziehen k\"{o}nnen.

Eine ver\"{a}nderliche Gr\"{o}sse, die f\"{u}r jeden Punkt~$O$
der Fl\"{a}che~$T$, allgemein zu reden, d.~h.\ ohne eine Ausnahme
in einzelnen Linien und Punkten\footnote{Diese Beschr\"{a}nkung
ist zwar nicht durch den Begriff einer Function an sich geboten,
aber um Infinitesimalrechnung auf sie anwenden zu k\"{o}nnen
erforderlich: eine Function, die in allen Punkten einer
Fl\"{a}che unstetig ist, wie z.~B.\ eine Function, die f\"{u}r
ein commensurables $x$ und ein commensurables $y$ den Werth~$1$,
sonst aber den Werth~$2$ hat, kann weder einer Differentiation,
noch einer Integration, also (unmittelbar) der
Infinitesimalrechnung \"{u}berhaupt nicht unterworfen werden.
Die f\"{u}r die Fl\"{a}che ~$T$ hier willk\"{u}rlich gemachte
Beschr\"{a}nking wird sich sp\"{a}ter (Art.~15) rechtfertigen.}
auszuschliessen, Einen bestimmten mit der Lage desselben stetig
sich \"{a}ndernden Werth annimmt, kann offenbar als eine Function
von $x$,~$y$, angesehen werden, und \"{u}berall, wo in der Folge
von Functionen von $x$,~$y$ die Rede sein wird, werden wir den
Begriff derselben auf diese Art festlegen.

Ehe wir uns jedoch zur Betrachtung solcher Functionen wenden,
schalten wir noch einige Er\"{o}rterungen \"{u}ber den
Zusammenhang einer Fl\"{a}che ein.  Wir beschr\"{a}nken uns dabei
auf solche Fl\"{a}chen, die sich nicht l\"{a}ngs einer Linie
spalten.

\medbreak

\centerline{6.}

\nobreak\medskip

Wir betrachten zwei Fl\"{a}chentheile als zusammenh\"{a}ngend
oder Einem St\"{u}cke angeh\"{o}rig, wenn sich von einem Punkte
des einen durch das Innere der Fl\"{a}che eine Linie nach einem
Punkte des andern ziehen l\"{a}sst, als getrennt, wenn diese
M\"{o}glichkeit nicht Statt findet.

Die Untersuchung des Zusammenhangs einer Fl\"{a}che beruht auf
ihrer Zerlegung durch Querschnitte, d.~h.\ Linien, welche von
einem Begrenzungspunkte das Innere einfach---keinen Punkt
mehrfach---bis zu einem Begrenzungspunkte durchschneiden.
Letzterer kann auch in dem zur Begrenzung hinzukommenen Theile,
also in einem fr\"{u}hern Punkte des Querschnitts, liegen.

Eine zusammenh\"{a}ngende Fl\"{a}che heisst, wenn sie durch jeden
Querschnitt in St\"{u}cke zerf\"{a}llt, eine einfach
zusammenh\"{a}ngende, andernfalls eine mehrfach
zusammenh\"{a}ngende.

\emph{Lehrsatz.~I}.  Eine einfach zusammenh\"{a}ngende
Fl\"{a}che~$A$ zerf\"{a}llt durch jeden Querschnitt $ab$ in zwei
einfach zusammenh\"{a}ngende St\"{u}cke.

Gesetzt, eins dieser St\"{u}cke w\"{u}rde durch einen Querschnitt
$cd$ nicht zerst\"{u}ckt, so erhielte man offenbar, je nachdem
keiner seiner Endpunkte oder der Endpunkt~$c$ oder beide
Endpunkte in $ab$ fielen, durch Herstellung der Verbindung
l\"{a}ngs der ganzen Linie $ab$ oder l\"{a}ngs des Theils $cb$
oder des Theils $cd$ derselben eine zusammenh\"{a}ngende
Fl\"{a}che, welche durch einen Querschnitt aus $A$ entst\"{a}nde,
gegen die Voraussetzung.

\emph{Lehrsatz.~II}.  Wenn eine Fl\"{a}che $T$ durch
$n_1$ Querschnitte\footnote{Unter einer Zerlegung durch mehrere
Querschnitte ist stets eine successive zu verstehen, d.~h.\ eine
solche, wo die durch einen Querschnitt \emph{entstandene} Fl\"{a}che
durch einen neuen Querschnitt weiter zerlegt wird.}
$q_1$ in ein System $T_1$ von $m_1$ einfach zusammenh\"{a}ngenden
Fl\"{a}chenst\"{u}cken und durch $n_2$ Querschnitte $q_2$ in ein
System $T_2$ von $m_2$ Fl\"{a}chenst\"{u}cken zerf\"{a}llt, so
kann $n_2 - m_2$ nicht $> n_1 - m_1$ sein.

Jede Linie $q_2$ bildet, wenn sie nicht ganz in das
Querschnittsystem $q_1$ f\"{a}llt, zugleich einen oder mehrere
Querschnitte $q_2'$ der Fl\"{a}che $T_1$.  Als Endpunkte der
Querschnitte $q_2'$ sind anzusehen:
\begin{description}
\item[\textmd{1)}]
die $2n_2$ Endpunkte der Querschnitte $q_2$, ausgenommen, wenn
ihre Enden mit einem Theil des Liniensystems $q_1$
zusammenfallen,
\item[\textmd{2)}]
jede mittlere Punkt eines Querschnitts $q_2$, in welchem er in
einen mittlern Punkt einer Linie $q_1$ eintritt, ausgenommen,
wenn er sich schon in einer andern Linie $q_1$ befindet,
d.~h.\ wenn ein Ende eines Querschnitts $q_1$ mit ihm
zusammenf\"{a}llt.
\end{description}

Bezeichnet nun $\mu$, wie oft Linien beider Systeme w\"{a}hrend
ihres Laufes zusammentreffen oder auseinandergehen (wo also ein
einzelner gemeinsamer Punkt doppelt zu rechnen ist), $\nu_1$, wie
oft ein Endst\"{u}ck der $q_1$ mit einem mittlern St\"{u}cke der
$q_2$, $\nu_2$, wie oft ein Endst\"{u}ck der $q_2$ mit einem
mittlern St\"{u}cke der $q_1$, endlich $\nu_3$, wie oft ein
Endst\"{u}ck der $q_1$ mit einem Endst\"{u}cke der $q_2$
zusammenf\"{a}llt, so liefert Nr.~1 $2n_2 - \nu_2 - \nu_3$,
Nr.~2 $\mu - \nu_1$ Endpunkte der Querschnitte $q_2'$; beide
F\"{a}lle zusammengenommen aber umfassen s\"{a}mmtliche Endpunkte
und jeden nur einmal, und die Anzahl dieser Querschnitte ist
daher
\[ \frac{2n_2 - \nu_2 - \nu_3 + \mu - \nu_1}{2}
   = n_2 + s.\]
Durch ganz \"{a}hnliche Schl\"{u}sse ergiebt sich die Anzahl der
Querschnitte $q_1'$ der Fl\"{a}che $T_2$, welche durch die Linien
$q_1$ gebildet werden,
\[ \frac{2n_1 - \nu_1 - \nu_3 + \mu - \nu_2}{2},\]
also $= n_1 + s$.  Die Fl\"{a}che $T_1$ wird nun offenbar durch
die $n_2 + s$ Querschnitte $q_2'$ in dieselbe Fl\"{a}che
verwandelt, in welche $T_2$ durch die $n_1 + s$ Querschnitte
$q_1'$ zerf\"{a}llt wird.  Es besteht aber $T_1$ aus $m_1$
einfach zusammenh\"{a}ngenden St\"{u}cken und zerf\"{a}llt daher
nach Satz~I durch $n_2 + s$ Querschnitte in $m_1+ n_2 + s$
Fl\"{a}chenst\"{u}cke; folglich m\"{u}sste, w\"{a}re
$m_2 < m_1 + n_2 - n_1$, die Zahl der Fl\"{a}chenst\"{u}cke $T_2$
durch $n_1 + s$ Querschnitte um mehr als $n_1 + s$ vermehrt
werden, was ungereimt ist.

Zufolge dieses Lehrsatzes ist, wenn die Anzahl der Querschnitte
unbestimmt durch $n$, die Anzahl der St\"{u}cke durch $m$
bezeichnet wird, $n - m$ f\"{u}r alle Zerlegungen einer
Fl\"{a}che in einfach zusammenh\"{a}ngende St\"{u}cke constant;
denn betrachten wir irgend zwei bestimmte Zerlegungen durch $n_1$
Querschnitte in $m_1$ St\"{u}cke und durch $n_2$ Querschnitte in
$m_2$ St\"{u}cke, so muss, wenn erstere einfach
zusammenh\"{a}ngend sind, $n_2 - m_2 \leq n_1 - m_1$, und wenn
letztere einfach zusammenh\"{a}ngend sind,
$n_1 - m_1 \leq n_2 - m_2$, also wenn Beides zutrifft,
$n_2 - m_2 = n_1 - m_1$ sein.

Diese Zahl kann f\"{u}glich mit dem Namen \glqq Ordnung des
Zusammenhangs\grqq\ einer Fl\"{a}che belegt werden; sie wird

durch jeden Querschnitt um $1$ erniedrigt---nach der
Definition---,

durch eine von einem innern Punkte das Innere einfach bis zu
einem Begrenzungspunkte oder einem fr\"{u}hern Schnittpunkte
durchschneidende Linie nicht ge\"{a}ndert und

durch einen innern allenthalben einfachen in zwei Punkten
endenden Schnitt um $1$ erh\"{o}ht,

weil erstere durch Einen, letzterer aber durch zwei Querschnitte
in Einen Querschnitt verwandelt werden kann.

Endlich wird die Ordnung des Zusammenhangs einer aus mehreren
St\"{u}cken bestehenden Fl\"{a}che erhalten, wenn man die
Ordnungen des Zusammenhangs dieser St\"{u}cke zu einander addirt.

Wir werden uns indess in der Folge meistens auf eine aus Einem
St\"{u}cke betstehende Fl\"{a}che beschr\"{a}nken, und uns
f\"{u}r ihren Zusammanhang der kunstloseren Bezeichnung eines
einfachen, zweifachen etc.\ bedienen, indem wir unter einer
$n \,$\-fach zusammenh\"{a}ngenden Fl\"{a}che eine solche
verstehen, die durch $n - 1$ Querschnitte in eine einfach
zusammenh\"{a}ngende zerlegbar ist.

In Bezug auf die Abh\"{a}ngigkeit des Zusammenhangs der
Begrenzung von dem Zusammenhang einer Fl\"{a}che erhallt leicht:

1)  Die Begrenzung einer einfach zusammenh\"{a}ngenden Fl\"{a}che
besteht noth\-wendig aus Einer in sich zur\"{u}cklaufenden Linie.

Best\"{a}nde die Begrenzung ausgetrennten St\"{u}cken, so
w\"{u}rde ein Querschnitt $q$, der einen Punkt eines St\"{u}cks
$a$ mit einem Punkte eines andern $b$ verb\"{a}nde, nur
zusammenh\"{a}ngende Fl\"{a}chentheile von einander scheiden, da
sich im Innern der Fl\"{a}che l\"{a}ngs $a$ eine Linie von der
einen Seite des Querschnitts $q$ an die entgegengesetzte
f\"{u}hren liesse; und folglich w\"{u}rde $q$ die Fl\"{a}che
nicht zerst\"{u}cken, gegen die Voraussetzung.

2)  Durch jeden Querschnitt wird die Anzahl der
Begrenzungsst\"{u}cke entweder um $1$ vermindert oder um $1$
vermehrt.

Ein Querschnitt $q$ verbindet entweder einen Punkt eines
Begrenzungsst\"{u}cks $a$ mit einem Punkte eines andern $b$,---in
diesem Falle bilden alle diese Linien zusammengenommen in der
Folge $a$,~$q$,~$b$,~$q$ ein einziges in sich zur\"{u}cklaufendes
St\"{u}ck der Begrenzung---

oder er verbindet zwei Punkte eines St\"{u}cks der
Begrenzung,---in diesem Falle zerf\"{a}llt dieses durch seine
beiden Endpunkte in zwei St\"{u}cke, deren jedes mit dem
Querschnitte zusammengenommen ein in sich zur\"{u}cklaufendes
Begrenzungsst\"{u}ck bildet---

oder endlich, er endet in einem seiner fr\"{u}heren Punkte und
kann betrachtet werden als zusammengesetzt aus einer in sich
zur\"{u}cklaufenden Linie $o$ und einer andern $l$, welche einen
Punkt von $o$ mit einem Punkte eines Begrenzungsst\"{u}cks $a$
verbindet,---in welchem Falle $o$ eines Theils, und
$a$,~$l$,$~o$,~$l$ andern Theils je ein in sich
zur\"{u}cklaufendes Begrenzungsst\"{u}ck bilden.

Es treten also entweder---im erstern Falle---an die Stelle zweier
Ein, oder---in den beiden letzteren F\"{a}llen---an die Stelle
Eines zwei Begrenzungsst\"{u}cke, woraus unser Satz folgt.

Die Anzahl der St\"{u}cke, aus welchen die Begrenzung eines
$n \,$\-fach zusammenh\"{a}ngenden Fl\"{a}chenst\"{u}cks besteht, ist
daher entweder $= n$ oder um eine gerade Zahl kleiner.

Hieraus ziehen wir noch das Corollar:

Wenn die Anzahl der Begrenzungsst\"{u}cke einer $n \,$\-fach
zusammenh\"{a}ngende Fl\"{a}che $= n$ ist, so zerf\"{a}llt
diese durch jeden \"{u}berall einfachen im Innern in sich
zur\"{u}cklaufenden Schnitt in zwei getrennte St\"{u}cke.

Denn die Ordnung des Zusammenhangs wird dadurch nicht
ge\"{a}ndert, die Anzahl der Begrenzungsst\"{u}cke um $2$
vermehrt; die Fl\"{a}che w\"{u}rde also, wenn sie eine
zusammenh\"{a}ngende w\"{a}re, einen $n \,$\-fachen Zusammenhang und
$n + 2$ Begrenzungsst\"{u}cke haben, was unm\"{o}glich ist.

\medbreak

\centerline{7.}

\nobreak\medskip

Sind $X$ und $Y$ zwei in allen Punkten der \"{u}ber $A$
ausgebreiteten Fl\"{a}che~$T$ stetige Functionen von $x$,~$y$, so
ist das \"{u}ber alle Elemente $dT$ dieser Fl\"{a}che ausgedehnte
Integral
\[ \int \left(
         \frac{\partial X}{\partial x}
       + \frac{\partial Y}{\partial y}
      \right) \, dT
   = - \int (X \cos \xi + Y \cos \eta) ds,\]
wenn in jedem Punkte der Begrenzung die Neigung einer auf sie
nach Innen gezogenen Normale gegen die $x$-Axe durch $\xi$, gegen
die $y$-Axe durch $\eta$ bezeichnet wird, und sich diese
Integration auf s\"{a}mmtliche Elemente $ds$ der Begrenzungslinie
erstreckt.

Um das Integral
$\displaystyle \int \frac{\partial X}{\partial x} \, dT$
zu transformiren, zerlegen wir den von der Fl\"{a}che~$T$
bedeckten Theil der Ebene $A$ durch ein System der $x$-Axe
paralleler Linien in Elementarstreifen, und zwar so, dass jeder
Windungspunkt der Fl\"{a}che~$T$ in eine dieser Linien f\"{a}llt.
Unter dieser Voraussetzung besteht der auf jeden derselben
fallende Theil von $T$ aus einem oder mehreren abgesondert
verlaufenden trapezf\"{o}rmigen St\"{u}cken.  Der Beitrag eines
unbestimmten dieser Fl\"{a}chenstreifen, welcher aus der $y$-Axe
das Element $dy$ ausscheidet, zu dem Werthe von
$\displaystyle \int \frac{\partial X}{\partial x} \, dT$
wird dann offenbar
$\displaystyle = dy \, \int \frac{\partial X}{\partial x} \, dx$,
wenn diese Integration durch diejenige oder diejenigen der
Fl\"{a}che~$T$ angeh\"{o}rigen geraden Linien ausgedehnt wird,
welche auf eine durch einen Punkt von $dy$ gehende Normale
fallen.  Sind nun die unteren Endpunkte derselben (d.~h.\ welchen
die kleisten Werthe von $x$ entsprechen)
$O_\prime, O_{\prime\prime}, O_{\prime\prime\prime},\ldots$,
die oberen
$O', O'', O''',\ldots$
und bezeichen wir mit
$X_\prime, X_{\prime\prime},\ldots$ $X', X'',\ldots$
die Werthe von $X$ in diesen Punkten, mit
$ds_\prime, ds_{\prime\prime},\ldots$ $ds', ds'',\ldots$
die entsprechenden von dem Fl\"{a}chenstreifen
aus der Begrenzung ausgeschiedenen Elemente, mit
$\xi_\prime, \xi_{\prime\prime},\ldots$ $\xi', \xi'',\ldots$
die Werthe von $\xi$ an diesen Elementen, so wird
\begin{eqnarray*}
\int \frac{\partial X}{\partial x} \, dx
   &=& - X_\prime - X_{\prime\prime} - X_{\prime\prime\prime} \ldots \\
   & & + X' + X'' + X''' \ldots .
\end{eqnarray*}
Die Winkel $\xi$ werden offenbar spitz an den unteren, stumpf and
den oberen Endpunkten, und es wird daher
\begin{eqnarray*}
dy &=&   \cos \xi_\prime \, ds_\prime
    =    \cos \xi_{\prime\prime} \, ds_{\prime\prime} \ldots \\
   &=& - \cos \xi' \, ds'
    =  - \cos \xi'' \, ds''\ldots .
\end{eqnarray*}
Durch Substitution dieser Werthe ergiebt sich
\[ dy \, \int \frac{\partial X}{\partial x} \, dx
   = - {\textstyle\sum} X \cos \xi \, ds,\]
wo sich die Summation auf alle Begrenzungselemente bezieht,
welche in der $y$-Axe $dy$ zur Projection haben.

Durch Integration \"{u}ber s\"{a}mmtliche in Betracht kommende
$dy$ werden offenbar s\"{a}mmtliche Elemente der Fl\"{a}che $T$
und s\"{a}mmtliche Elemente der Begrenzung ersch\"{o}pft, und man
erh\"{a}lt daher, in diesem Umfange genommen,
\[ \int \frac{\partial X}{\partial x} \, dT
   = - \int X \cos \xi \, ds.\]
Durch ganz \"{a}hnliche Schl\"{u}sse findet man
\[ \int \frac{\partial Y}{\partial y} \, dT
   = - \int Y \cos \eta \, ds \]
und folglich
\[ \int \left(
         \frac{\partial X}{\partial x}
       + \frac{\partial Y}{\partial y}
      \right) \, dT
   = - \int (X \cos \xi + Y \cos \eta) ds, \quad\mbox{w.~z.~b.~w.} \]

\medbreak

\centerline{8.}

\nobreak\medskip

Bezeichnen wir in der Begrenzungslinie, von einem festen
Anfangspunkte auf in einer bestimmten sp\"{a}ter festzusetzenden
Richtung gerechnet, die L\"{a}nge derselben bis zu einem
unbestimmten Punkte $O_0$, durch $s$, und in der in diesem Punkte
$O_0$ errichteten Normalen die Entfernung eines unbestimmten
Punktes $O$ von demselben und zwar nach Innen zu als positiv
betrachtet durch $p$, so k\"{o}nnen offenbar die Werthe von $x$
und $y$ im Punkte $O$ als Functionen von $s$ und $p$ angesehen
werden, und es werden dann in den Punkten der Begrenzungslinie
die partiellen Differentialquotienten
\[ \frac{\partial x}{\partial p} = \cos \xi,\quad
   \frac{\partial y}{\partial p} = \cos \eta,\quad
   \frac{\partial x}{\partial s} = \pm \cos \eta,\quad
   \frac{\partial y}{\partial s} = \mp \cos \xi,\]
wo die oberen Zeichen gelten, wenn die Richtung, in welcher die
Gr\"{o}sse $s$ als wachsend betrachtet wird, mit $p$ einen
gleichen Winkel einschliesst, wie die $x$-Axe mit der $y$-Axe,
wenn einen entgegengesetzten, die unteren.  Wir werden diese
Richtung un allen Theilen der Begrenzung so annehmen, dass
\[ \frac{\partial x}{\partial s}
      = \frac{\partial y}{\partial p}
\quad\mbox{und folglich}\quad
   \frac{\partial y}{\partial s}
      = - \frac{\partial x}{\partial p} \]
ist, was die Allgemeinheit unserer Resultate im Wesentlichen
nicht beeintr\"{a}chtigt.

Offenbar k\"{o}nnen wir diese Bestimmungen auch auf Linien im
Innern von $T$ ausdehnen; nur haben wir hier zur Bestimmung der
Vorzeichen von $dp$ und $ds$, wenn deren gegenseitige
Abh\"{a}ngigkeit wie dort festgesetzt wird, noch eine Angabe
hinzuzuf\"{u}gen, welche entweder das Vorzeichen von $dp$ oder
von $ds$ festsetzt; und zwar werden wir bei einer in sich
zur\"{u}cklaufenden Linie angeben, von welchem der durch sie
geschiedenen Fl\"{a}chentheile sie als Begrenzung gelten solle,
wodurch das Vorzeichen von $dp$ bestimmt wird, bei einer nicht in
sich zur\"{u}cklaufenden aber ihren Anfangspunkt, d.~h. den
Endpunkt, wo $s$ den kleinsten Werth annimmt.

Die Einf\"{u}hrung der f\"{u}r $\cos \xi$ und $\cos \eta$
erhaltenen Werthe in die im vorigen Art.\ bewiesene Gleichung
giebt, in demselben Umfange wie dort genommen,
\[ \int  \left(
            \frac{\partial X}{\partial x}
          + \frac{\partial Y}{\partial y}
         \right) \, dT
  = - \int \left(
            X \frac{\partial x}{\partial p}
          + Y \frac{\partial y}{\partial p}
         \right) \, ds
  = \int \left(
            X \frac{\partial y}{\partial s}
          - Y \frac{\partial x}{\partial s}
         \right) \, ds.\]

\medbreak

\centerline{9.}

\nobreak\medskip

Durch Anwendung des Satzes am Schlusse des vorigen Art.\ auf den
Fall, wo in allen Theilen der Fl\"{a}che
\[ \frac{\partial X}{\partial x} + \frac{\partial Y}{\partial y}
   = 0 \]
ist, erhalten wir folgende S\"{a}tze:

1.  Sind $X$ und $Y$ zwei in allen Punkten von $T$ endliche und
stetige und der Gleichung
\[ \frac{\partial X}{\partial x} + \frac{\partial Y}{\partial y}
   = 0 \]
gen\"{u}gende Functionen, so ist, durch die ganze Begrenzung von
$T$ ausgedehnt,
\[ \int  \left(
            X \frac{\partial x}{\partial p}
          + Y \frac{\partial y}{\partial p}
         \right) \, ds
    = 0.\]

Denkt man sich eine beliebige \"{u}ber $A$ ausgestreckte
Fl\"{a}che $T_1$ in zwei St\"{u}cke $T_2$ und $T_3$ auf
beliebige Art zerf\"{a}llt, so kann das Integral
\[ \int  \left(
            X \frac{\partial x}{\partial p}
          + Y \frac{\partial y}{\partial p}
         \right) \, ds \]
in Bezug auf die Begrenzung von $T_2$ betrachtet werden als die
Differenz der Integrale in Bezug auf die Begrenzung von $T_1$ und
in Bezug auf die Begrenzung von $T_3$, indem, wo $T_3$ sich bis
zur Begrenzung von $T_1$ erstreckt, beide Integrale sich
aufheben, alle \"{u}brigen Elemente aber einem Elemente der
Begrenzung von $T_2$ entsprechen.

Mittelst dieser Umformung ergiebt sich aus I.:

II.  Der Werth des Integrals
\[ \int  \left(
            X \frac{\partial x}{\partial p}
          + Y \frac{\partial y}{\partial p}
         \right) \, ds,\]
durch die ganze Begrenzung einer \"{u}ber $A$ ausgebreiteten
Fl\"{a}che erstreckt, bleibt bei beliebiger Erweiterung oder
Verengerung derselben constant, wenn nur dadurch keine
Fl\"{a}chentheile ein- oder austreten, innerhalb welcher die
Voraussetzungen des Satzes~I.\ nicht erf\"{u}llt sind.

Wenn die Functionen $X$, $Y$ zwar in jedem Theile der Fl\"{a}che
$T$ der vorgeschriebenen Differentialgleichung gen\"{u}gen, aber
in einzelnen Linien oder Punkten mit einer Unstetigkeit behaftet
sind, so kann man jede solche Linie und jeden solchen Punkt mit
einem beliebig kleinen Fl\"{a}chentheil als H\"{u}lle umgeben und
erh\"{a}lt dann durch Anwendung des Satzes~II.:

III.  Das Integral
\[ \int  \left(
            X \frac{\partial x}{\partial p}
          + Y \frac{\partial y}{\partial p}
         \right) \, ds \]
in Bezug auf die ganze Begrenzung von $T$ ist gleich der Summe der
Integrale
\[ \int  \left(
            X \frac{\partial x}{\partial p}
          + Y \frac{\partial y}{\partial p}
         \right) \, ds \]
in Bezug auf die Umgrenzungen aller Unstetigkeitsstellen, und zwar
beh\"{a}lt in Bezug auf jede einzelne dieser Stellen das Integral
denselben Werth, in wie enge Grenzen man sie auch einschliessen
m\"{o}ge.

Dieser Werth ist f\"{u}r einen blossen Unstetigkeitspunkt
nothwendig gleich Null, wenn mit der Entfernung $\varrho$ des
Punktes $O$ von demselben zugleich $\varrho X$ und $\varrho Y$
unendlich klein werden; denn f\"{u}hrt man in Bezug auf einen
solchen Punkt als Anfangspunkt und eine beliebige Anfangsrichtung
Polarcoordinaten $\varrho$, $\varphi$ ein und w\"{a}hlt zur
Umgrenzung einen um denselben mit dem Radius $\varrho$
beschriebenen Kreis, so wird das auf ihn bez\"{u}gliche Integral
durch
\[ \int\limits_{0}^{2\pi} \left(
            X \frac{\partial x}{\partial p}
          + Y \frac{\partial y}{\partial p}
         \right) \varrho \, d\varphi \]
ausgedr\"{u}ckt und kann folglich nicht einen von Null
verschiedenen Werth $\kappa$ haben, weil, was auch $\kappa$ sei,
$\varrho$ immer so klein angenommen werden kann, dass abgesehen
vom Zeichen
$\displaystyle \left(
            X \frac{\partial x}{\partial p}
          + Y \frac{\partial y}{\partial p}
         \right) \varrho$
f\"{u}r jeden Werth von $\varphi$ kleiner als
$\displaystyle \frac{\kappa}{2\pi}$
und folglich
\[ \int\limits_{0}^{2\pi} \left(
            X \frac{\partial x}{\partial p}
          + Y \frac{\partial y}{\partial p}
         \right) \varrho \, d\varphi
   < \kappa \]
wird.

IV.  Ist in einer einfach zusammenh\"{a}ngenden \"{u}ber $A$
ausgebreiteten Fl\"{a}che f\"{u}r jeden Fl\"{a}chentheil das
durch dessen ganze Begrenzung erstreckte Integral
\[ \int  \left(
            X \frac{\partial x}{\partial p}
          + Y \frac{\partial y}{\partial p}
         \right) \, ds \]
oder
\[ \int  \left(
            Y \frac{\partial x}{\partial s}
          - X \frac{\partial y}{\partial s}
         \right) \, ds
   = 0,\]
so erh\"{a}lt f\"{u}r irgend zwei feste Punkte $O_0$ und $O$ dies
Integral in Bezug auf alle von $O_0$ in derselben nach $O$
gehende Linien denselben Werth.

Je zwei die Punkte $O_0$ und $O$ verbindende Linien $s_1$ und
$s_2$ bilden zusammengenommen eine in sich zur\"{u}cklaufende
Linie $s_3$. Diese Linie besitzt entweder selbst die Eigenschaft,
keinen Punkt mehrfach zu durchschneiden, oder man kann sie in
mehrere allenthalben einfache in sich zur\"{u}cklaufende Linien
zerlegen, indem man von einem beliebigen Punkte aus dieselbe
durchlaufend jedesmal, wenn man zu einem fr\"{u}hern Punkte
zur\"{u}ckgelangt, den inzwischen durchlaufenen Theil ausscheidet
und den folgenden als unmittelbare Fortsetzung des vorhergehenden
betrachtet.  Jede solche Linie aber zerlegt die Fl\"{a}che in
eine einfach und eine zweifach zusammenh\"{a}ngende; sie bildet
daher nothwendig von Einem dieser St\"{u}cke die ganze
Begrenzung, und das durch sie erstreckte Integral
\[ \int  \left(
            Y \frac{\partial x}{\partial s}
          - X \frac{\partial y}{\partial s}
         \right) \, ds \]
wird also der Voraussetzung nach $= 0$.  Dasselbe gilt folglich
auch von dem durch die ganze Linie $s_3$ erstreckten Integrale,
wenn die Gr\"{o}sse $s$ \"{u}berall in derselben Richtung als
wachsend betrachtet wird; es m\"{u}ssen daher die durch die
Linien $s_1$ und $s_2$ erstreckten Integrale, wenn diese
Richtung unge\"{a}ndert bleibt, d.~h.\ in einer derselben von
$O_0$ nach $O$ und in der andern von $O$ nach $O_0$ geht,
einander aufheben, also, wenn sie in letzterer ge\"{a}ndert wird,
gleich werden.

Hat man nun irgend eine beliebige Fl\"{a}che~$T$, in welcher,
allgemein zu reden,
\[ \frac{\partial X}{\partial x} + \frac{\partial Y}{\partial y}
   = 0 \]
ist, so schliesse man zun\"{a}chst, wenn n\"{o}thig, die
Unstetigkeitsstellen aus, so dass im \"{u}brigen
Fl\"{a}chenst\"{u}cke f\"{u}r jeden Fl\"{a}chentheil
\[ \int  \left(
            Y \frac{\partial x}{\partial s}
          - X \frac{\partial y}{\partial s}
         \right) \, ds
   = 0 \]
ist, und zerlege dieses durch Querschnitte in eine einfach
zusammenh\"{a}ngende Fl\"{a}che $T^*$.  F\"{u}r jede im Innern
von $T^*$ von einem Punkte $O_0$ nach einem andern $O$ gehende
Linie hat dann unser Integral demselben Werth; dieser Werth,
f\"{u}r den zur Abk\"{u}rzung die Bezeichnung
\[ \int\limits_{O_0}^O  \left(
            Y \frac{\partial x}{\partial s}
          - X \frac{\partial y}{\partial s}
         \right) \, ds \]
gestattet sein m\"{o}ge, ist daher $O_0$ als fest, $O$ als
beweglich gedacht, f\"{u}r jede Lage von $O$, abgesehen vom Laufe
der Verbindungslinie ein bestimmter, und kann folglich als
Function von $x$,~$y$ betrachtet werden.  Die Aenderung dieser
Function wird f\"{u}r eine Verr\"{u}ckung von $O$ l\"{a}ngs eines
beliebigen Linienelements $ds$ durch
\[ \left(
            Y \frac{\partial x}{\partial s}
          - X \frac{\partial y}{\partial s}
         \right) \, ds \]
ausgedr\"{u}ckt, ist in $T^*$ \"{u}berall stetig und l\"{a}ngs
eines Querschnitts von $T$ zu beiden Seiten gleich;

V.  das Integral
\[ Z = \int\limits_{O_0}^O \left(
            Y \frac{\partial x}{\partial s}
          - X \frac{\partial y}{\partial s}
         \right) \, ds \]
bildet daher, $O_0$ als fest gedacht, eine Function von $x$,~$y$,
welche in $T^*$ \"{u}berall sich stetig, beim Ueberschreiten der
Querschnitte von $T$ aber um eine l\"{a}ngs derselben von einem
Zweigpunkte zum andern constante Gr\"{o}sse \"{a}ndert, und von
welcher der partielle Differentialquotient
\[ \frac{\partial Z}{\partial x} = Y,\quad
   \frac{\partial Z}{\partial y} = - X \]
ist.

Die Aenderungen beim Ueberschreiten der Querschnitte sind von
einer der Zahl der Querschnitte gleichen Anzahl von einander
unabh\"{a}ngiger Gr\"{o}ssen abh\"{a}ngig; denn wenn man das
Querschnittsystem r\"{u}ckw\"{a}rts---die sp\"{a}teren Theile
zuerst---durchl\"{a}uft, so ist diese Aenderung \"{u}berall
bestimmt, wenn ihr Werth beim Beginn jedes Querschnitts gegeben
wird; letztere Werthe aber sind von einander unabh\"{a}ngig.

\medbreak

\centerline{10.}

\nobreak\medskip

Setzt man f\"{u}r die bisher durch $X$ bezeichnete Function
\[ u \frac{\partial u'}{\partial x}
   - u' \frac{\partial u}{\partial x} \]
und
\[ u \frac{\partial u'}{\partial y}
   - u' \frac{\partial u}{\partial y} \]
f\"{u}r $Y$, so wird
\[ \frac{\partial X}{\partial x} + \frac{\partial Y}{\partial y}
  =   u  \left(
            \frac{\partial^2 u'}{\partial x^2}
          + \frac{\partial^2 u'}{\partial y^2}
         \right)
    - u' \left(
            \frac{\partial^2 u }{\partial x^2}
          + \frac{\partial^2 u }{\partial y^2}
         \right),\]
wenn also die Functionen $u$ und $u'$ den Gleichungen
\[ \frac{\partial^2 u }{\partial x^2} + \frac{\partial^2 u }{\partial y^2}
      = 0,\quad
   \frac{\partial^2 u'}{\partial x^2} + \frac{\partial^2 u'}{\partial y^2}
      = 0 \]
gen\"{u}gen, so wird
\[ \frac{\partial X}{\partial x} + \frac{\partial Y}{\partial y}
   = 0,\]
und es finden auf den Ausdruck
\[ \int  \left(
            X \frac{\partial x}{\partial p}
          + Y \frac{\partial y}{\partial p}
         \right) \, ds,\]
welcher
\[ = \int \left(
            u  \frac{\partial u'}{\partial p}
          - u' \frac{\partial u }{\partial p}
         \right) \, ds \]
wird, die S\"{a}tze des vorigen Art.\ Anwendung.

Machen wir nun in Bezug auf die Function~$u$ die Voraussetzung,
dass sie nebst ihren ersten Differentialquotienten etwaige
Unstetigkeiten jedenfalls nicht l\"{a}ngs einer Linie erleidet,
und f\"{u}r jeden Unstetigkeitspunkt zugleich mit der
Entfernung~$\varrho$ des Punktes~$O$ von demselben
$\displaystyle \varrho \frac{\partial u}{\partial x}$
und
$\displaystyle \varrho \frac{\partial u}{\partial y}$
unendlich klein werden, so k\"{o}nnen die Unstetigkeiten von $u$
in Folge der Bermerkung zu III. des vorigen Art.\ ganz
unber\"{u}cksichtigt bleiben.

Denn alsdann kann man in jeder von einem Unstetigkeitspunkte
ausgehenden geraden Linie einen Werth $R$ von $\varrho$ so
annehmen, dass
\[ \varrho \frac{\partial u}{\partial \varrho}
   =   \varrho \frac{\partial u}{\partial x}
               \frac{\partial x}{\partial \varrho}
     + \varrho \frac{\partial u}{\partial y}
               \frac{\partial y}{\partial \varrho} \]
unterhalb desselben immer endlich bleibt, und bezeichnet $U$ den
Werth von $u$ f\"{u}r $\varrho = R$, $M$ abgesehen von Zeichen
den gr\"{o}ssten Werth der Function
$\displaystyle \varrho \frac{\partial u}{\partial \varrho}$
in jenem Intervall, so wird, in derselben Bedeutung genommen,
stets
$u - U < M (\log \varrho - \log R)$
sein, folglich $\varrho (u - U)$ und also auch $\varrho u$ mit
$\varrho$ zugleich unendlich klein werden; dasselbe gilt aber der
Voraussetzung nach von
$\displaystyle \varrho \frac{\partial u}{\partial x}$
und
$\displaystyle \varrho \frac{\partial u}{\partial y}$
und folglich wenn $u'$ keiner Unstetigkeit unterliegt, auch von
\[ \varrho \left( u \frac{\partial u'}{\partial x}
   - u' \frac{\partial u}{\partial x} \right)
      \quad\mbox{und}\quad
   \varrho \left( u \frac{\partial u'}{\partial y}
   - u' \frac{\partial u}{\partial y} \right);\]
der im vorigen Art.\ er\"{o}rterte Fall tritt hier also ein.

Wir nehmen nun ferner an, dass die den Ort des Punktes $O$
bildende Fl\"{a}che $T$ allenthalben einfach \"{u}ber $A$
ausgebreitet sei, und denken uns in derselben einen beliebigen
festen Punkt $O_0$, wo $u$, $x$,~$y$ die Werthe
$u_0$, $x_0$,~$y_0$ erhalten.  Die Gr\"{o}sse
\[ {\textstyle\frac{1}{2}} \log
         ((x - x_0)^2 + (y - y_0)^2) = \log r,\]
als Function von $x$,~$y$ betrachtet, hat alsdann die
Eigenschaft, dass
\[    \frac{\partial^2 \log r}{\partial x^2}
    + \frac{\partial^2 \log r}{\partial y^2}
   = 0 \]
wird, und ist nur f\"{u}r $x = x_0$,~$y = y_0$, also in unserm
Falle nur f\"{u}r Einen Punkt der Fl\"{a}che $T$ mit einer
Unstetigkeit behaftet.

Es wird daher nach Art.~9, III., wenn wir $\log r$ f\"{u}r $u'$
setzen,
\[ \int  \left(
            u  \frac{\partial \log r}{\partial p}
          - \log r \frac{\partial u }{\partial p}
         \right) \, ds \]
in Bezug auf die ganze Begrenzung von $T$ gleich diesem Integrale
in Bezug auf eine beliebige Umgrenzung der Punktes $O_0$ und
also, wenn wir dazu die Peripherie eines Kreises, wo $r$ einen
constanten Werth hat, w\"{a}hlen und von einem ihrer Punkte in
einer beliebigen festen Richtung den Bogen bis $O$ in Theilen des
Halbmessers durch $\varphi$ bezeichnen, gleich
\[ - \int u  \frac{\partial \log r}{\partial r} r \, d \varphi
      - \log r \int \frac{\partial u }{\partial p} \, ds,\]
oder da
\[ \int \frac{\partial u }{\partial p} \, ds = 0 \]
ist,
\[ = - \int_0^{2\pi} u \, d \varphi,\]
welcher Werth, wenn $u$ im Punkte $O_0$ stetig ist, f\"{u}r ein
unendlich kleines $r$ in $-u_0 \, 2\pi$ \"{u}bergeht.

Unter den in Bezug auf $u$ und $T$ gemachten Voraussetzungen
haben wir daher f\"{u}r einen beliebigen Punkt $O_0$ im Innern der
Fl\"{a}che, in welchem $u$ stetig ist,
\[ u_0 = \frac{1}{2\pi} \int \left(
            \log r \frac{\partial u }{\partial p}
          - u  \frac{\partial \log r}{\partial p}
         \right) \, ds \]
in Bezug auf die ganze Begrenzung derselben und
\[ = \frac{1}{2\pi} \int_0^{2\pi} u \, d \varphi \]
in Bezug auf einen um $O_0$ beschriebenen Kreis.  Aus dem ersten
dieser Ausdr\"{u}cke ziehen wir folgenden

\emph{Lehrsatz}.  Wenn eine Function $u$ innerhalb einer die
Ebene $A$ allenthalben einfach bedeckenden Fl\"{a}che $T$,
allgemein zu reden, der Differentialgleichung
\[    \frac{\partial^2 u }{\partial x^2}
    + \frac{\partial^2 u }{\partial y^2}
   = 0 \]
gen\"{u}gt und zwar so, dass
\begin{description}

\item[\textmd{1)}]
die Punkte, in welchen diese Differentialgleichung nicht
erf\"{u}llt ist, keinen Fl\"{a}chentheil,

\item[\textmd{2)}]
die Punkte, in welchen $u$,
$\displaystyle \frac{\partial u}{\partial x}$,
$\displaystyle \frac{\partial u}{\partial y}$
unstetig werden, keine Linie stetig erf\"{u}llen,

\item[\textmd{3)}]
f\"{u}r jeden Unstetigkeitspunkt zugleich mit der Entfernung
$\varrho$ des Punktes $O$ von demselben die Gr\"{o}ssen
$\displaystyle \varrho \frac{\partial u}{\partial x}$,
$\displaystyle \varrho \frac{\partial u}{\partial y}$
unendlich klein werden und

\item[\textmd{4)}]
bei $u$ eine durch Ab\"{a}nderung ihres Werthes in einzelnen
Punkten hebbare Unstetigkeit ausgeschlossen ist,
\end{description}

\noindent
so ist sie nothwendig nebst allen ihren Differentialquotienten
f\"{u}r alle Punkte im Innern dieser Fl\"{a}che endlich und
stetig.

In der That, betrachten wir den Punkt $O_0$ als beweglich, so
\"{a}ndern sich in dem Ausdrucke
\[ \int \left(
            \log r \frac{\partial u }{\partial p}
          - u  \frac{\partial \log r}{\partial p}
         \right) \, ds \]
nur die Werthe
$\log r$,
$\displaystyle \varrho \frac{\partial \log r}{\partial x}$,
$\displaystyle \varrho \frac{\partial \log r}{\partial y}$.
Diese Gr\"{o}ssen aber sind f\"{u}r jedes Element der Begrenzung,
so lange $O_0$ in Innern von $T$ bleibt, nebst allen ihren
Differentialquotienten endliche und stetige Functionen von
$x_0$,~$y_0$, da die Differentialquotienten durch gebrochene
rationale Function dieser Gr\"{o}ssen ausgedr\"{u}ckt werden, die
nur Potenzen von $r$ in Nenner enthalten.  Dasselbe gilt daher
auch f\"{u}r den Werth unsres Integrals und folglich f\"{u}r die
Function $u_0$.  Denn diese k\"{o}nnte unter den fr\"{u}heren
Voraussetzungen nur in einzelnen Punkten, indem sie unstetig
w\"{u}rde, einen davon verschiedenen Werth haben, welche
M\"{o}glichkeit durch die Voraussetzung 4) unsers Lehrsatzes
wegf\"{a}llt.

\medbreak

\centerline{11.}

\nobreak\medskip

Under denselben Voraussetzungen in Bezug auf $u$ und $T$, wie am
Schlusse des vorigen Art., haben wir folgende S\"{a}tze:

I.  Wenn l\"{a}ngs einer Linie $u = 0$ und
$\displaystyle \frac{\partial u}{\partial p} = 0$
ist, so ist $u$ \"{u}berall $= 0$.

Wir beweisen zun\"{a}chst, dass eine Linie $\lambda$, wo $u = 0$ und
$\displaystyle \frac{\partial u}{\partial p} = 0$
ist, nicht die Begrenzung eines Fl\"{a}chentheils $a$, wo $u$
positiv ist, bilden k\"{o}nne.

Gesetzt, dies f\"{a}nde statt, so scheide man aus $a$ ein Stuck
aus, welches eines Theils durch $\lambda$, andern Theils durch
eine Kreislinie begrenzt wird und den Mittelpunkt $O_0$ dieses
Kreises nicht enth\"{a}lt, welche Construction allemal
m\"{o}glich ist.  Man hat dann, wenn man die Polarcoordinaten von
$O$ in Bezug auf $O_0$ durch $r$, $\varphi$ bezeichnet, durch die
ganze Begrenzung dieses St\"{u}cks ausgedehnt
\[ \int \log r \frac{\partial u }{\partial p} \, ds
          - \int u  \frac{\partial \log r}{\partial p} \, ds
   = 0,\]
also in Folge der Annahme auch f\"{u}r den ganzen ihr
angeh\"{o}rigen Kreisbogen
\[ \int u \, d \varphi
      + \log r \int \frac{\partial u }{\partial p} \, ds
   = 0,\]
oder da
\[ \int \frac{\partial u}{\partial p} \, ds = 0 \]
ist,
\[ \int u \, d \varphi = 0,\]
was mit der Voraussetzung, dass $u$ im Innern von $a$ positiv
sei, unvertr\"{a}glich ist.

Auf \"{a}hnliche Art wird bewiesen, dass die Gleichungen $u = 0$
und
$\displaystyle \frac{\partial u}{\partial p} = 0$
nicht in einem Begrenzungstheile eines Fl\"{a}chenst\"{u}cks~$b$,
wo $u$ negativ ist, stattfinden k\"{o}nnen.

Wenn nun in der Fl\"{a}che $T$ in einer Linie $u = 0$ und
$\displaystyle \frac{\partial u}{\partial p} = 0$
ist und in irgend einem Theile derselben $u$ von Null verschieden
w\"{a}re, so m\"{u}sste ein solcher Fl\"{a}chentheil offenbar
entweder durch diese Linie selbst oder durch einen
Fl\"{a}chentheil, wo $u = 0$ w\"{a}re, also jedenfalls durch eine
Linie wo $u$ und
$\displaystyle \frac{\partial u}{\partial p} = 0$
w\"{a}re, begrenzt werden, was nothwendig auf eine der vorhin
widerlegten Annahmen f\"{u}hrt.

II.  Wenn der Werth von $u$ und
$\displaystyle \frac{\partial u}{\partial p}$
l\"{a}ngs einer Linie gegeben ist, so ist $u$ dadurch in allen
Theilen von $T$ bestimmt.

Sind $u_1$ und $u_2$ irgend zwei bestimmte Functionen, welche den
der Function $u$ auferlegten Bedingungen gen\"{u}gen, so gilt
dies auch, wie sich durch Substitution in diese Bedingungen
sofort ergiebt, f\"{u}r ihre Differenz $u_1 - u_2$.  Stimmten nun
$u_1$ und $u_2$ l\"{a}ngs einer Linie nebst ihren ersten
Differentialquotienten nach $p$ \"{u}berein, in einem andern
Fl\"{a}chentheile aber nicht, so w\"{u}rden l\"{a}ngs dieser
Linie $u_1 - u_2 = 0$ und
$\displaystyle \frac{\partial (u_1 - u_2)}{\partial p} = 0$
sein, ohne \"{u}berall $= 0$ zu sein, dem Satze~I. zuwider.

III.  Die Punkte im Innern von $T$, wo $u$ einen constanten Werth
hat, bilden, wenn $u$ nicht \"{u}berall constant ist, nothwendig
Linien, welche Fl\"{a}chentheile, wo $u$ gr\"{o}sser ist, von
Fl\"{a}chentheilen, wo $u$ kleiner ist, scheiden.

Dieser Satz ist aus folgenden zusammengesetzt:

$u$ kann nicht in einem Punkte im Innern von $T$ ein Minimum oder
ein Maximum haben;

$u$ kann nicht \emph{nur} in einem Theile der Fl\"{a}che constant
sein;

die Linien, in denen $u = a$ ist, k\"{o}nnen nicht beiderseits
Fl\"{a}chentheile begrenzen, wo $u - a$ dasselbe Zeichen hat;

S\"{a}tze, deren Gegentheil, wie leicht zu sehen, allemal eine
Verletzung der im vorigen Art.\ bewiesenen Gleichung
\[ u_0 = \frac{1}{2\pi} \int_0^{2\pi} u \, d \varphi \]
oder
\[ \int_0^{2\pi} (u - u_0) \, d \varphi  = 0 \]
herbeif\"{u}hren m\"{u}sste und folglich unm\"{o}glich ist.

\medbreak

\centerline{12.}

\nobreak\medskip

Wir wenden uns jetzt zur\"{u}ck zur Betrachtung einer
ver\"{a}nderlichen complexen Gr\"{o}sse $w = u + vi$, welche,
allgemein zu reden (d.~h.\ ohne eine Ausnahme in einzelnen Linien
und Punkten auszuschliessen), f\"{u}r jeden Punkt $O$ der
Fl\"{a}che $T$ Einen bestimmten mit der Lage desselben stetig und
den Gleichungen
\[ \frac{\partial u}{\partial x}
      = \frac{\partial v}{\partial y},\quad
   \frac{\partial u}{\partial y}
      = - \frac{\partial v}{\partial x} \]
gem\"{a}ss sich \"{a}ndernden Werth hat, und bezeichnen diese
Eigenschaft von $w$ nach dem fr\"{u}her Festgestellten dadurch,
dass wir $w$ eine Function von $z = x + yi$ nennen.  Zur
Vereinfachung des Folgenden setzen wir dabei im Voraus fest, dass
bei einer Function von $z$ eine durch Ab\"{a}nderung ihres
Werthes in einem einzelnen Punkte hebbare Unstetigkeit nicht
vorkommen solle.

Der Fl\"{a}che $T$ wird vorerst ein einfacher Zusammenhang und eine
allenthalben einfache Ausbreitung \"{u}ber die Ebene $A$
beigelegt.

\emph{Lehrsatz}.  Wenn eine Function $w$ von $z$ eine
Unterbrechung der Stetigkeit jedenfalls nicht l\"{a}ngs einer
Linie erleidet und ferner f\"{u}r jeden beliebigen Punkt $O'$ der
Fl\"{a}che, wo $z = z'$ sei, $w(z - z')$ mit unendlicher
Ann\"{a}herung des Punktes~$O$ unendlich klein wird, so ist sie
nothwendig nebst allen ihren Differentialquotienten in allen
Punkten im Innern der Fl\"{a}che endlich und stetig.

Die \"{u}ber die Ver\"{a}nderungen der Gr\"{o}sse $w$ gemachten
Voraussetzungen zerfallen, wenn
$z - z' = \varrho e^{\varphi i}$
gesetzt wird, f\"{u}r $u$ und $v$ in die folgenden:
\[ \mbox{1)}\quad
   \frac{\partial u}{\partial x} - \frac{\partial v}{\partial y}
      = 0 \]
und
\[ \mbox{2)}\quad
   \frac{\partial u}{\partial y} + \frac{\partial v}{\partial x}
      = 0 \]
f\"{u}r jeden Theil der Fl\"{a}che $T$;
\begin{description}
\item[\textmd{3)}]
die Functionen $u$ und $v$ sind nicht l\"{a}ngs einer Linie
unstetig;
\item[\textmd{4)}]
f\"{u}r jeden Punkt $O'$ werden mit der Entfernung $\varrho$ des
Punktes $O$ von demselben $\varrho u$ und $\varrho v$ unendlich
klein;
\item[\textmd{5)}]
f\"{u}r die Functionen $u$ und $v$ sind Unstetigkeiten, die durch
Ab\"{a}nderung ihres Werthes in einzelnen Punkten gehoben werden
k\"{o}nnten, ausgeschlossen.
\end{description}

In Folge der Voraussetzungen 2), 3), 4) ist f\"{u}r jeden Theil
der Fl\"{a}che $T$ das \"{u}ber dessen ganze Begrenzung
ausgedehnte Integral
\[ \int  \left(
            u \frac{\partial x}{\partial s}
          - v \frac{\partial y}{\partial s}
         \right) \, ds \]
nach Art.~9, III. $= 0$ und das Integral
\[ \int\limits_{O_0}^O
         \left(
            u \frac{\partial x}{\partial s}
          - v \frac{\partial y}{\partial s}
         \right) \, ds \]
erh\"{a}lt daher (nach Art.~9, IV.) durch jede von $O_0$ nach $O$
gehende Linie erstreckt denselben Werth und bildet, $O_0$ als
fest gedacht, eine bis auf einzelne Punkte nothwendig stetige
Function $U$ von $x$,~$y$, von welcher (und zwar nach 5) in jedem
Punkte) der Differentialquotient
$\displaystyle \frac{\partial U}{\partial x} = u$
und
$\displaystyle \frac{\partial U}{\partial y} = - v$
ist.  Durch Substitution dieser Werthe f\"{u}r $u$ und $v$ aber
gehen die Voraussetzungen 1), 3), 4), in die Bedingungen des
Lehrsatzes am Schlusse des Art.~10 \"{u}ber.  Die Function $U$ ist
daher nebst allen ihren Differentialquotienten in allen Punkten
von $T$ endlich und stetig und dasselbe gilt folglich auch von
der complexen Function
$\displaystyle w
   =  \frac{\partial U}{\partial x} - \frac{\partial U}{\partial y} i$
und ihren nach $z$ genommenen Differentialquotienten.

\medbreak

\centerline{13.}

\nobreak\medskip

Es soll jetzt untersucht werden, was eintritt, wenn wir unter
Beibehaltung der sonstigen Voraussetzungen des Art.~12 annehmen,
dass f\"{u}r einen bestimmten Punkt $O'$ im Innern der Fl\"{a}che
$(z - z') w = \varrho e^{\varphi i} w$
bei unendlicher Ann\"{a}herung des Punktes $O$ nicht mehr
unendlich klein wird.  In diesem Falle wird also $w$ bei
unendlicher Ann\"{a}herung des Punktes $O$ an $O'$ unendlich
gross, und wir nehmen an, dass, wenn die Gr\"{o}sse $w$ nicht mit
$\displaystyle \frac{1}{\varrho}$
von gleicher Ordnung bleibt, d.~h.\ der Quotient beider sich
einer endlichen Grenze n\"{a}hert, wenigstens die Ordnungen
beider Gr\"{o}ssen in einem endlichen Verh\"{a}ltnisse zu
einander stehen, so dass sich eine Potenz von $\varrho$ angeben
l\"{a}sst, deren Product in $w$ f\"{u}r ein unendlich kleines
$\varrho$ entweder unendlich klein wird oder endlich bleibt.  Ist
$\mu$ der Exponent einer solchen Potenz und $n$ die n\"{a}chst
gr\"{o}ssere ganze Zahl, so wird die Gr\"{o}sse
$(z - z')^n w = \varrho^n e^{n\varphi i} w$
mit $\varrho$ unendlich klein, und es ist daher
$(z - z')^{n-1} w$ eine Function von $z$ (da
$\displaystyle \frac{d(z - z')^{n-1} w}{dz}$
von $dz$ unabh\"{a}ngig ist), welche in diesem Theile der
Fl\"{a}che den Voraussetzungen des Art.~12 gen\"{u}gt und
folglich im Punkte $O'$ endlich und stetig ist.  Bezeichnen wir
ihren Werth im Punkte $O'$ mit $a_{n-1}$, so ist
$(z - z')^{n-1} w - a_{n-1}$
eine Function, die in diesem Punkte stetig und $= 0$ ist und
folglich mit $\varrho$ unendlich klein wird, woraus man nach
Artikel~12 schliesst, dass
$\displaystyle (z - z')^{n-2} w - \frac{a_{n-1}}{z - z'}$
eine im Punkte $O'$ stetige Function ist.  Durch Forsetzung
dieses Verfahrens wird offenbar $w$ mittelst Subtraction eines
Ausdruckes von der Form
\[ \frac{a_1}{z - z'}
      + \frac{a_2}{(z - z')^2} + \cdots
      + \frac{a_{n-1}}{(z - z')^{n-1}} \]
in eine Function verwandelt, welche im Punkte $O'$ endlich und
stetig bleibt.

Wenn daher unter den Voraussetzungen des Art.~12 die Aenderung
eintritt, dass bei unendlicher Ann\"{a}herung von $O$ an einen
Punkt $O'$ im Innern der Fl\"{a}che $T$ die Function $w$
unendlich gross wird, so ist die Ordnung dieses unendlich Grossen
(eine im verkehrten Verh\"{a}ltnisse der Entfernung wachsende
Gr\"{o}sse als ein unendlich Grosses erster Ordnung betrachtet)
wenn sie endlich ist, nothwendig eine ganze Zahl; und ist diese
Zahl $= m$, so kann die Function $w$ durch Hinzuf\"{u}gung einer
Function, welche $2m$ willk\"{u}rliche Constanten enth\"{a}lt,
in eine in diesem Punkte $O'$ stetige verwandelt werden.

\medskip

{\footnotesize
Anm.  Wir  betrachten eine Function als Eine willk\"{u}rliche
Constante enthaltend, wenn die m\"{o}glichen Arten, sie zu
bestimmen, ein stetiges Gebiet von Einer Dimension umfassen.\par}

\medbreak

\centerline{14.}

\nobreak\medskip

Die im Art.~12 und 13 in Bezug auf die Fl\"{a}che~$T$ gemachten
Beschr\"{a}nkungen sind f\"{u}r die G\"{u}ltigkeit der gewonnenen
Resultate nicht wesentlich.  Offenbar kann man jeden Punkt im
Innern einer beliebigen Fl\"{a}che mit einem St\"{u}cke derselben
umgeben, welches die dort vorausgesetzten Eigenschaften besitzt,
mit alleiniger Ausnahme des Falles, wo dieser Punkt ein
Windungspunkt der Fl\"{a}che ist.

Um diesen Fall zu untersuchen, denken wir uns die Fl\"{a}che~$T$
oder ein beliebiges St\"{u}ck derselben, welches einen
Windungspunkt $(n - 1) \,$\-ter Ordnung $O'$, wo $z = z' = x' + y' i$
sei, enth\"{a}lt, mittelst der Function
$\displaystyle \zeta = (z - z')^{\frac{1}{n}}$
auf einer andern Ebene $\Lambda$ abgebildet, d.~h.\ wir denken
uns den Werth der Function $\zeta = \xi + \eta i$ im Punkte~$O$
durch einen Punkt $\Theta$, dessen rechtwinklige Coordinaten
$\xi$,~$\eta$ sind, in dieser Ebene vertreten, und betrachten
$\Theta$ als Bild des Punktes $O$.  Auf diesem Wege erh\"{a}lt
man als Abbildung dieses Theils der Fl\"{a}che~$T$ eine
zusammenh\"{a}ngende \"{u}ber $\Lambda$ ausgebreitete Fl\"{a}che,
die im Punkte $\Theta'$, dem Bilde des Punktes $O'$, keinen
Windungspunkt hat, wie sogleich gezeigt werden soll.

Zur Fixirung der Vorstellungen denke man sich um den Punkt~$O$ in
der Ebene~$A$ mit dem Halbmesser~$R$ einen Kreis beschrieben und
parallel mit der $x$-Axe einen Durchmesser gezogen, wo also
$z - z'$ reelle Werthe annehmen wird.  Das durch diesen Kreis
ausgeschiedene den Windingspunkt umgebende St\"{u}ck der
Fl\"{a}che $T$ wird dann zu beiden Seiten des Durchmessers in
$n$, wenn $R$ hinreichend klein gew\"{a}hlt wird, abgesondert
verlaufende halbkreisf\"{o}rmige Fl\"{a}chenst\"{u}cke zerfallen.
Wir bezeichnen auf derjenigen Seite des Durchmessers, wo $y - y'$
positiv ist, diese Fl\"{a}chenst\"{u}cke durch
$a_1, a_2 \,\ldots\, a_n$, auf der entgegengesetzten Seite durch
$a_1', a_2' \,\ldots\, a_n'$, und nehmen an, dass f\"{u}r
negative Werthe von $z - z'$ $a_1, a_2 \,\ldots\, a_n$ der Reihe
nach mit $a_1', a_2' \,\ldots\, a_n'$, f\"{u}r positive dagegen
mit $a_n', a_1' \,\ldots\, a_{n-1}'$ verbunden seien, so dass ein
den Punkt $O'$ (im erforderlichen Sinne) umkreisender Punkt der
Reihe nach die Fl\"{a}chen
$a_1, a_1', a_2, a_2' \,\ldots\, a_n, a_n'$
durchl\"{a}uft und durch $a_n'$ wieder in $a_1$
zur\"{u}ckgelangt, welche Annahme offenbar gestattet ist.
F\"{u}hren wir nun f\"{u}r beide Ebenen Poloarcoordinaten ein,
indem wir $z - z' = \varrho e^{\varphi i}$,
$\zeta = \sigma e^{\psi i}$ setzen, und w\"{a}hlen zur Abbildung
des Fl\"{a}chenst\"{u}cks $a_1$ denjenigen Werth von
$\displaystyle (z - z')^{\frac{1}{n}}
     = \varrho^{\frac{1}{n}} e^{\frac{\varphi}{n}i}$,
welchen letzterer Ausdruck unter der Annahme
$0 \leq \varphi \leq \pi$
erh\"{a}lt, so wird f\"{u}r alle Punkte von $a_1$
$\displaystyle \sigma \leq R^{\frac{1}{n}}$
und
$\displaystyle 0 \leq \psi \leq \frac{\pi}{n}$;
die Bilder derselben in der Ebene $\Lambda$ fallen also
s\"{a}mmtlich in einen von $\psi = 0$ bis
$\displaystyle \psi = \frac{\pi}{n}$
sich erstreckenden Sector eines um $\Theta'$ mit dem Radius
$\displaystyle R^{\frac{1}{n}}$
beschriebenen Kreises, und zwar entspricht jedem Punkte von $a_1$
Ein zugleich mit demselben stetig fortr\"{u}ckender Punkt dieses
Sectors und umgekehrt, woraus folgt, dass die Abbildung der
Fl\"{a}che $a_1$ eine zusammenh\"{a}ngende einfach \"{u}ber
diesen Sector ausgebreitete Fl\"{a}che ist.  Auf \"{a}hnliche Art
erh\"{a}lt man f\"{u}r die Fl\"{a}che $a_1'$ als Abbildung einen
von $\displaystyle \psi = \frac{\pi}{n}$
bis $\displaystyle \psi = \frac{2\pi}{n}$,
for $a_2$ einen
von $\displaystyle \psi = \frac{2\pi}{n}$
bis $\displaystyle \psi = \frac{3\pi}{n}$,
endlich f\"{u}r $a_n'$ einen
von $\displaystyle \psi = \frac{2n - 1}{n}\pi$
bis $\displaystyle \psi = 2\pi$
sich erstreckenden Sector, wenn man $\varphi$ f\"{u}r jeden Punkt
dieser Fl\"{a}chen der Reihe nach zwischen $\pi$ und $2\pi$,
$2\pi$ und $3\pi$ \dots\ $(2n - 1)\pi$ und $2n\pi$ w\"{a}hlt, was
immer und nur auf eine Weise m\"{o}glich ist.  Diese Sectoren
schliessen sich aber in derselben Folge an einander, wie die
Fl\"{a}chen $a$ und $a'$, und zwar so, dass den hier
zusammenstossenden Punkten auch dort zusammenstossende Punkte
entsprechen; sie k\"{o}nnen daher zu einer zusammenh\"{a}ngenden
Abbildung eines den Punkt $O'$ einschliessenden St\"{u}ckes der
Fl\"{a}che~$T$ zusammengef\"{u}gt werden, und diese Abbilding ist
offenbar eine \"{u}ber die Ebene $\Lambda$ einfach ausgebreitete
Fl\"{a}che.

Eine ver\"{a}nderliche Gr\"{o}sse, die f\"{u}r jeden Punkt~$O$
einen bestimmten Werth hat, hat dies auch f\"{u}r jeden Punkt
$\Theta$ und umgekehrt, da jedem $O$ nur ein $\Theta$ und jedem
$\Theta$ nur ein $O$ entspricht; ist sie ferner eine Function von
$z$, so ist sie dies auch von $\zeta$, indem, wenn
$\displaystyle \frac{dw}{dz}$
von $dz$, auch
$\displaystyle \frac{dw}{d\zeta}$
von $d\zeta$ unabh\"{a}ngig ist, und umgekehrt.  Es ergiebt sich
hieraus, dass auf alle Functionen $w$ von $z$ auch im
Windungspunkte $O'$ die S\"{a}tze der Art.~12 und 13 angewandt
werden k\"{o}nnen, wenn man sie als Functionen von
$\displaystyle (z - z')^{\frac{1}{n}}$
betrachtet.  Dies liefert folgenden Satz:

Wenn eine Function $w$ von $z$ bei unendlicher Ann\"{a}herung von
$O$ an einen Windungspunkt $(n - 1) \,$\-ter Ordnung $O'$ unendlich
wird, so ist dieses unendlich Grosse nothwendig von gleicher
Ordnung mit einer Potenz der Entfernung, deren Exponent ein
Vielfaches von
$\displaystyle \frac{1}{n}$
ist, und kann, wenn dieser Exponent
$\displaystyle = - \frac{m}{n}$
ist, durch Hinzuf\"{u}gung eines Ausdrucks von der Form
\[    \frac{a_1}{(z - z')^{\frac{1}{n}}}
    + \frac{a_2}{(z - z')^{\frac{2}{n}}} + \cdots
    + \frac{a_m}{(z - z')^{\frac{m}{n}}},\]
wo $a_1, a_2 \,\ldots\, a_m$ willk\"{u}rliche complexe
Gr\"{o}ssen sind, in eine im Punkte $O'$ stetige verwandelt
werden.

Dieser Satz enh\"{a}lt als Corollar, dass die Function~$w$ im
Punkte~$O'$ stetig ist, wenn
$\displaystyle (z - z')^{\frac{1}{n}} w$
bei unendlicher Ann\"{a}herung des Punktes $O$ an $O'$ unendlich
klein wird.

\medbreak

\centerline{15.}

\nobreak\medskip

Denken wir uns jetzt eine Function von $z$, welche f\"{u}r jeden
Punkt~$O$ der beliebig \"{u}ber $A$ ausgebreiteten Fl\"{a}che~$T$
einen bestimmten Werth hat und nicht \"{u}berall constant ist,
geometrisch dargestellt, so dass ihr Werth $w = u + vi$ im
Punkte~$O$ durch einen Punkt~$Q$ der Ebene~$B$ vertreten wird,
dessen rechtwinklige Coordinaten $u$,~$v$ sind, so ergiebt sich
Folgendes:

I.  Die Gesammtheit der Punkte~$Q$ kann betrachtet werden als
eine Fl\"{a}che $S$ bildend, in welcher jedem Punkte Ein
bestimmter mit ihm stetig in $T$ fortr\"{u}ckender Punkt $O$
entspricht.

Um dieses zu beweisen, ist offenbar nur der Nachweis
erforderlich, dass die Lage des Punktes~$Q$ mit der des
Punktes~$O$ sich allemal (und zwar, allgemein zu reden, stetig)
\"{a}ndert.  Dieser ist in dem Satze enthalten:

Eine Function $w = u + vi$ von $z$ kann nicht l\"{a}ngs einer
Linie constant sein, wenn sie nicht \"{u}berall constant ist.

Beweis: H\"{a}tte $w$ l\"{a}ngs einer Linie einen constanten
Werth $a + bi$ so w\"{a}ren $u - a$ und
$\displaystyle \frac{\partial (u - a)}{\partial p}$,
welches
$\displaystyle = - \frac{\partial v}{\partial s}$,
f\"{u}r diese Linie und
\[    \frac{\partial^2 (u - a)}{\partial x^2}
    + \frac{\partial^2 (u - a)}{\partial y^2} \]
\"{u}berall $= 0$; es m\"{u}sste also nach Art.~11, I.\ $u - a$
und folglich, da
\[ \frac{\partial u}{\partial x}
      = \frac{\partial v}{\partial y},\quad
   \frac{\partial u}{\partial y}
      = - \frac{\partial v}{\partial x},\]
auch $v - b$ \"{u}berall $= 0$ sein, gegen die Voraussetzung.

II.  In Folge der in I. gemachten Voraussetzung kann zwischen den
Theilen von $S$ nicht ein Zusammenhang Statt finden ohne einen
Zusammenhang der entsprechenden Theile von $T$; umgekehrt kann
\"{u}berall, wo in $T$ Zusammenhang Statt findet und $w$ stetig
ist, der Fl\"{a}che $S$ ein entsprechender Zusammenhang beigelegt
werden.

Dieses vorausgesetzt entspricht die Begrenzung von $S$
einestheils der Begrenzung von $T$, anderntheils den
Unstetigkeitsstellen; ihre inneren Theile aber sind, einzelne
Punkte ausgenommen, \"{u}berall schlicht \"{u}ber $B$
ausgebreitet, d.~h.\ es findet nirgends eine Spaltung in auf
einander liegende Theile und nirgends eine Umfaltung Statt.

Ersters k\"{o}nnte, da $T$ \"{u}berall einen entsprechenden
Zusammenhang besitzt, offenbar nur eintreten, wenn in $T$ eine
Spaltung vork\"{a}me---der Annahme zuwider---; Letzteres soll
sogleich bewiesen werden.

Wir beweisen zuv\"{o}rderst, dass ein Punkt~$Q'$, wo
$\displaystyle \frac{dw}{dz}$
endlich ist, nicht in einer Falte der Fl\"{a}che $S$ liegen kann.

In der That, umgeben wir den Punkt~$O'$, welcher $Q'$ entspricht,
mit einem St\"{u}cke der Fl\"{a}che~$T$ von beliebiger Gestalt
und unbestimmten Dimensionen, so m\"{u}ssen (nach Art.~3) die
Dimensionen desselben stets so klein angenommen werden
k\"{o}nnen, dass die Gestalt des entsprechenden Theils von $S$
beliebig wenig abweicht, und folglich so klein, dass die
Begrenzung desselben aus der Ebene~$B$ ein $Q'$ einschliessendes
St\"{u}ck ausscheidet.  Dies aber ist unm\"{o}glich, wenn $Q'$
in einer Falte der Fl\"{a}che $S$ liegt.

Nun kann
$\displaystyle \frac{dw}{dz}$,
als Function von $z$, nach I.\ nur in einzelnen Punkten $= 0$,
und, da $w$ in den in Betracht kommenden Punkten von $T$ stetig
ist, nur in den Windungspunkten dieser Fl\"{a}che unendlich
werden; folglich etc.\ w.~z.~b.~w.

III.  Die Fl\"{a}che $S$ ist folglich eine Fl\"{a}che, f\"{u}r
welche die im Art.~5 f\"{u}r $T$ gemachten Voraussetzungen
zutreffen; und in dieser Fl\"{a}che hat f\"{u}r jeden Punkt~$Q$
die unbestimmte Gr\"{o}sse $z$ Einen bestimmten Werth, welcher
sich mit der Lage von $Q$ stetig und so \"{a}ndert, dass
$\displaystyle \frac{dz}{dw}$
von der Richtung der Orts\"{a}nderung unabh\"{a}ngig ist.  Es
bildet daher in dem fr\"{u}her festgelegten Sinne $z$ eine
stetige Function der ver\"{a}nderlichen complexen Gr\"{o}sse $w$
f\"{u}r das durch $S$ dargestellte Gr\"{o}ssengebiet.

Hieraus folgt ferner:

Sind $O'$ und $Q'$ zwei entsprechende innere Punkte der
Fl\"{a}chen $T$ und $S$ und in denselben $z = z'$, $w = w'$, so
n\"{a}hert sich, wenn keiner von ihnen ein Windungspunkt ist, bei
unendlicher Ann\"{a}herung von $O$ an $O'$
$\displaystyle \frac{w - w'}{z - z'}$
einer endlichen Grenze, und die Abbilding ist daselbst eine in
den kleinsten Theilen \"{a}hnliche; wenn aber $Q'$ ein
Windungspunkt $(n - 1) \,$\-ter, $O'$ ein Windungspunkt $(m - 1) \,$\-ter
Ordnung ist, so n\"{a}hert sich
$\displaystyle \frac{(w - w')^{\frac{1}{n}}}{(z - z')^{\frac{1}{m}}}$
bei unendlicher Ann\"{a}herung von $O$ an $O'$ einer endlichen
Grenze, und f\"{u}r die anstossenden Fl\"{a}chentheile findet
eine Abbildungsart Statt, die sich leicht aus Art.~14 ergiebt.

\nobreak\medskip

\asterisktriangle

\medbreak

\centerline{16.}

\nobreak\medskip

\emph{Lehrsatz}.  Sind $\alpha$ und $\beta$ zwei beliebige
Functionen von $x$,~$y$, f\"{u}r welche das Integral
\[ \int
      \left[
         \left(
            \frac{\partial \alpha}{\partial x}
          - \frac{\partial \beta }{\partial y}
         \right)^2
       + \left(
            \frac{\partial \alpha}{\partial y}
          + \frac{\partial \beta }{\partial x}
         \right)^2
      \right] \, dT \]
durch alle Theile der beliebig \"{u}ber $A$ ausgebreiteten
Fl\"{a}che $T$ ausgedehnt einen endlichen Werth hat, so
erh\"{a}lt das Integral bei Aenderung von $\alpha$ um stetige
oder doch nur in einzelnen Punkten unstetige Functionen, die am
Rande $= 0$ sind, immer f\"{u}r eine dieser Functionen einen
Minimumwerth und, wenn man durch Ab\"{a}nderung in einzelnen
Punkten hebbare Unstetigkeiten ausschliesst, nur f\"{u}r Eine.

Wir bezeichnen durch $\lambda$ eine unbestimmte stetige oder doch
nur in einzelnen Punkten unstetige Function, welche am Rande
$= 0$ ist und f\"{u}r welche das Integral
\[ L= \int
      \left(
         \left( \frac{\partial \lambda}{\partial x} \right)^2
       + \left( \frac{\partial \lambda}{\partial y} \right)^2
      \right) \, dT \]
\"{u}ber die ganze Fl\"{a}che ausgedehnt einen endlichen Werth
erh\"{a}lt, durch $\omega$ eine unbestimmte der Functionen
$\alpha + \lambda$, endlich das \"{u}ber die ganze Fl\"{a}che
erstreckte Integral
\[ \int
      \left[
         \left(
            \frac{\partial \omega}{\partial x}
          - \frac{\partial \beta }{\partial y}
         \right)^2
       + \left(
            \frac{\partial \omega}{\partial y}
          + \frac{\partial \beta }{\partial x}
         \right)^2
      \right] \, dT \]
durch $\Omega$.  Die Gesammtheit der Functionen $\lambda$ bildet
ein zusammenh\"{a}ngendes in sich abgeschossenes Gebiet, indem
jede dieser Functionen stetig in jede andere \"{u}bergehen, sich
aber nicht einer l\"{a}ngs einer Linie unstetigen unendlich
ann\"{a}hen kann, ohne dass $L$ unendlich wird (Art.~17); f\"{u}r
jedes $\lambda$ erh\"{a}lt nun, $\omega = \alpha + \lambda$
gesetzt, $\Omega$ einen endlichen Werth, der mit $L$ zugleich
unendlich wird, sich mit der Gestalt von $\lambda$ stetig
\"{a}ndert, aber nie unter Null herabsinken kann; folglich hat
$\Omega$ wenigstens f\"{u}r Eine Gestalt der Function $\omega$
ein Minimum.

Um den zweiten Theil unseres Satzes zu beweisen, sei $u$ eine der
Functionen $\omega$, welche $\Omega$ einen Minimumwerth ertheilt,
$h$ eine unbestimmte in der ganzen Fl\"{a}che constante
Gr\"{o}sse, so dass $u + h \lambda$ den der Function $\omega$
vorgeschriebenen Bedingungen gen\"{u}gt.  Der Werth von $\Omega$
f\"{u}r $\omega = u + h \lambda$, welcher
\begin{eqnarray*}
   &=& \int
      \left[
         \left(
            \frac{\partial u     }{\partial x}
          - \frac{\partial \beta }{\partial y}
         \right)^2
       + \left(
            \frac{\partial u     }{\partial y}
          + \frac{\partial \beta }{\partial x}
         \right)^2
      \right] \, dT \\
   & &+ 2 h \int
      \left[
         \left(
            \frac{\partial u     }{\partial x}
          - \frac{\partial \beta }{\partial y}
         \right) \frac{\partial \lambda}{\partial x}
       + \left(
            \frac{\partial u     }{\partial y}
          + \frac{\partial \beta }{\partial x}
         \right) \frac{\partial \lambda}{\partial y}
      \right] \, dT \\
   & &+ h^2 \int
      \left(
         \left( \frac{\partial \lambda}{\partial x} \right)^2
       + \left( \frac{\partial \lambda}{\partial y} \right)^2
      \right) \, dT \\
   &=& M + 2 N h + L h^2
\end{eqnarray*}
wird, muss alsdann f\"{u}r jedes $\lambda$ (nach dem Begriffe des
Minimums) gr\"{o}sser als $M$ werden, sobald $h$ nur hinreichend
klein genommen ist.  Dies erfordert aber, dass f\"{u}r jedes
$\lambda$ $N = 0$ sei; denn andernfalls w\"{u}rde
\[ 2 N h + L h^2 = L h^2 \left( 1 + \frac{2N}{Lh} \right) \]
negativ werden, wenn $h$ dem $N$ entgegengesetzt und abgesehen
vom Zeichen
$\displaystyle < \frac{2N}{L}$
angenommen w\"{u}rde.  Der Werth von $\Omega$ f\"{u}r
$\omega = u + \lambda$, in welcher Form offenbar alle
m\"{o}glichen Werthe von $\omega$ enthalten sind, wird daher
$= M + L$, und folglich kann, da $L$ wesentlich positiv ist,
$\Omega$ f\"{u}r keine Gestalt der Function $\omega$ einen
kleinern Werth erhalten, als f\"{u}r $\omega = u$.

Findet nun f\"{u}r eine andere $u'$ der Functionen $\omega$ ein
Minimumwerth $M'$ von $\Omega$ statt, so muss von diesem offenbar
dasselbe gelten, man hat also $M' \leq M$ und $M \leq M'$,
folglich $M = M'$.  Bringt man aber $u'$ auf die Form
$u + \lambda'$, so erh\"{a}lt man f\"{u}r $M'$ den Ausdruck
$M + L'$, wenn $L'$ den Werth von $L$ f\"{u}r $\lambda =
\lambda'$ bezeichnet, und die Gleichung $M = M'$ giebt $L' = 0$.
Dies ist nur m\"{o}glich, wenn in allen Fl\"{a}chentheilen
\[ \frac{\partial \lambda'}{\partial x} = 0,\quad
   \frac{\partial \lambda'}{\partial y} = 0 \]
ist, und es hat daher, so weit $\lambda'$ stetig ist, diese
Function nothwending einen constanten und folglich, da sie am
Rande $= 0$ und nicht l\"{a}ngs einer Linie unstetig ist,
h\"{o}chstens in einzelnen Punkten einen von Null verschiedenen
Werth.  Zwei der Functionen $\omega$, welche $\Omega$ einem
Minimumwerth ertheilen, k\"{o}nnen also nur in einzelnen Punkten
von einander verschieden sein, und wenn in der Function~$u$ alle
durch Ab\"{a}nderung in einzelnen Punkten hebbaren Unstetigkeiten
beseitigt werden, ist diese vollkommen bestimmt.

\medbreak

\centerline{17.}

\nobreak\medskip

Es soll jetzt der Beweis nachgeliefert werden, dass $\lambda$
unbeschadet der Endlichkeit von $L$ sich nicht einer l\"{a}ngs
einer Linie unstetigen Function $\gamma$ unendlich ann\"{a}hern
k\"{o}nne, d.~h.\ wird die Function $\lambda$ der Bedingung
unterworfen, ausserhalb eines die Unstetigkeitslinie
einschliessenden Fl\"{a}chentheils $T'$ mit $\gamma$
\"{u}bereinzustimmen, so kann $T'$ stets so klein angenommen
werden, das $L$ gr\"{o}sser als eine beliebig gegebene Gr\"{o}sse
$C$ werden muss.

Wir bezeichnen, $s$ und $p$ in Bezug auf die Unstetigkeitslinie
in der gewohnten Bedeutung genommen, f\"{u}r ein unbestimmtes $s$
die Kr\"{u}mmung, eine auf der Seite der positiven $p$ convexe
als positiv betrachtet, durch $\kappa$, den Werth von $p$ an der
Grenze von $T'$ auf der positiven Seite durch $p_1$, auf der
negativen Seite durch $p_2$ und die entsprechenden Werthe von
$\gamma$ durch $\gamma_1$ und $\gamma_2$.  Betrachten wir nun
irgend einen stetig gekr\"{u}mmten Theil dieser Linie, so liefert
der zwischen den Normalen in den Endpunkten enthaltene Theil von
$T'$, wenn er sich nicht bis zu den Kr\"{u}mmungsmittelpunkten
erstreckt, zu $L$ den Beitrag
\[ \int ds \int_{p_2}^{p_1} dp \, (1 - \kappa p)
         \left[
            \left( \frac{\partial \lambda}{\partial p} \right)^2
          + \left( \frac{\partial \lambda}{\partial s} \right)^2
               \frac{1}{(1 - \kappa p)^2}
         \right];\]
den kleinste Werth des Ausdrucks
\[ \int_{p_2}^{p_1}
            \left( \frac{\partial \lambda}{\partial p} \right)^2
               (1 - \kappa p) \, dp \]
bei den festen Grenzwerthen $\gamma_1$ und $\gamma_2$ von
$\lambda$ findet sich aber nach bekannten Regeln
\[ = \frac{(\gamma_1 - \gamma_2)^2 \kappa}{\log (1 - \kappa p_2)
                  - \log (1 - \kappa p_1)},\]
und folglich wird jener Beitrag nothwendig, wie auch $\lambda$
innerhalb $T'$ angenommen werden m\"{o}ge,
\[ > \int \frac{(\gamma_1 - \gamma_2)^2 \kappa \, ds}{\log (1 - \kappa p_2)
                  - \log (1 - \kappa p_1)}.\]
Die Function $\gamma$ w\"{a}re f\"{u}r $p = 0$ stetig, wenn der
gr\"{o}sste Werth, den $(\gamma_1 - \gamma_2)^2$ f\"{u}r
$\pi_1 > p_1 > 0$ und $\pi_2 < p_2 < 0$ erhalten kann, mit
$\pi_1 - \pi_2$ unendlich klein w\"{u}rde; wir k\"{o}nnen
folglich f\"{u}r jeden Werth von $s$ eine endliche Gr\"{o}sse $m$
so annehmen, dass, wie klein auch $\pi_1 - \pi_2$ angenommen
werden m\"{o}ge, stets innerhalb der durch $\pi_1 > p_1 \geq 0$
und $\pi_2 < p_2 \leq 0$ (wo die Gleichheiten sich gegenseitig
ausschliessen) ausgedr\"{u}ckten Grenzen Werthe von $p_1$ und
$p_2$ enthalten sind, f\"{u}r welche
$(\gamma_1 - \gamma_2)^2 > m$
wird.  Nehmen wir ferner unter dem fr\"{u}heren
Beschr\"{a}nkungen eine Gestalt von $T'$ beliebig an, indem wir
$p_1$ und $p_2$ bestimmte Werthe $P_1$ und $P_2$ beilegen, und
bezeichnen den Werth des durch in Betracht gezogenen Theil der
Unstetigkeitslinie ansgedehnten Integrals
\[ \int \frac{m \kappa \, ds}{\log (1 - \kappa P_2)
                  - \log (1 - \kappa P_1)} \]
durch $a$, so k\"{o}nnen wir offenbar
\[   \int \frac{(\gamma_1 - \gamma_2)^2 \kappa \, ds}{\log (1 - \kappa p_2)
                  - \log (1 - \kappa p_1)}
   > C \]
machen, indem wir $p_1$ und $p_2$ f\"{u}r jeden Werth von $s$ so
annehmen, dass den Ungleichheiten
\[ p_1 < \frac{1 - (1 - \kappa P_1)^{\frac{a}{C}}}{\kappa},\quad
   p_2 > \frac{1 - (1 - \kappa P_2)^{\frac{a}{C}}}{\kappa},\quad
   \hbox{und}\quad (\gamma_1 - \gamma_2)^2 > m \]
gen\"{u}gt wird.  Dies aber hat zur Folge, dass, wie auch
$\lambda$ innerhalb $T'$ angenommen werden m\"{o}ge, der aus dem
in Betracht gezogenen St\"{u}cke von $T'$ stammende Theil von $L$
und folglich um so mehr $L$ selbst $> C$ wird, w.~z.~b.~w.

\medbreak

\centerline{18.}

\nobreak\medskip

Nach Art.~16 haben wir f\"{u}r die dort festgelegte Function~$u$
und f\"{u}r irgend eine der Functionen~$\lambda$
\[ N = \int 
      \left[
         \left(
            \frac{\partial u     }{\partial x}
          - \frac{\partial \beta }{\partial y}
         \right) \frac{\partial \lambda}{\partial x}
       + \left(
            \frac{\partial u     }{\partial y}
          + \frac{\partial \beta }{\partial x}
         \right) \frac{\partial \lambda}{\partial y}
      \right] \, dT,\]
durch die ganze Fl\"{a}che~$T$ ausgedehnt, $= 0$.  Aus dieser
Gleichung sollen jetzt weitere Schl\"{u}sse gezogen werden.

Scheidet man aus der Fl\"{a}che~$T$ ein die Unstetigkeitsstellen
von $u$, $\beta$, $\lambda$ einschliessendes St\"{u}ck $T'$ aus,
so findet sich der von dem \"{u}brigen St\"{u}cke~$T''$
herr\"{u}hrende Theil von $N$ mit H\"{u}lfe der Art.~7, 8, wenn
man
$\displaystyle
   \left(
      \frac{\partial u     }{\partial x}
    - \frac{\partial \beta }{\partial y}
   \right) \lambda$
f\"{u}r $X$ und
$\displaystyle
   \left(
      \frac{\partial u     }{\partial y}
    + \frac{\partial \beta }{\partial x}
   \right) \lambda$
f\"{u}r $Y$ setzt,
\[ = - \int \lambda
            \left(
               \frac{\partial^2 u}{\partial x^2}
             + \frac{\partial^2 u}{\partial y^2}
            \right) \, dT
     - \int
            \left(
               \frac{\partial u     }{\partial p}
             + \frac{\partial \beta }{\partial s}
            \right) \lambda \, ds.\]
In Folge der der Function~$\lambda$ auferlegten Grenzbedingung
wird der auf das mit $T$ gemeinschaftliche Begrenzungsst\"{u}ck
von $T''$ bez\"{u}gliche Theil von
\[     \int \left(
               \frac{\partial u     }{\partial p}
             + \frac{\partial \beta }{\partial s}
            \right) \lambda \, ds \]
gleich $0$, so dass $N$ betrachtet werden kann als
zusammengesetzt aus dem Integral
\[   - \int \lambda
            \left(
               \frac{\partial^2 u}{\partial x^2}
             + \frac{\partial^2 u}{\partial y^2}
            \right) \, dT \]
in Bezug auf $T''$ und
\[ \int
      \left[
         \left(
            \frac{\partial u     }{\partial x}
          - \frac{\partial \beta }{\partial y}
         \right) \frac{\partial \lambda}{\partial x}
       + \left(
            \frac{\partial u     }{\partial y}
          + \frac{\partial \beta }{\partial x}
         \right) \frac{\partial \lambda}{\partial y}
      \right] \, dT
   + \int
            \left(
               \frac{\partial u     }{\partial p}
             + \frac{\partial \beta }{\partial s}
            \right) \lambda \, ds \]
in Bezug auf $T'$.

Offenbar w\"{u}rde nun, wenn
$\displaystyle 
   \frac{\partial^2 u}{\partial x^2}
 + \frac{\partial^2 u}{\partial y^2}$
in irgend einem Theile der Fl\"{a}che~$T$ von $0$ verscheiden
w\"{a}re, $N$ ebenfalls einen von $0$ verschiedenen Werth
erhalten, sobald man $\lambda$, was frei steht, innerhalb $T'$
gleich $0$ und innerhalb $T''$ so w\"{a}hlte, dass
$\displaystyle \lambda
   \left(
      \frac{\partial^2 u}{\partial x^2}
    + \frac{\partial^2 u}{\partial y^2}
   \right)$
\"{u}berall dasselbe Zeichen h\"{a}tte.  Ist aber
$\displaystyle 
   \frac{\partial^2 u}{\partial x^2}
 + \frac{\partial^2 u}{\partial y^2}$
in allen Theilen von $T$ $= 0$, so verschwindet der von $T''$
herr\"{u}hrende Bestandtheil von $N$ f\"{u}r jedes $\lambda$, und
die Bedingung $N = 0$ ergiebt dann, dass die auf die
Unstetigkeitsstellen bez\"{u}glichen Bestandtheile $= 0$ werden.

F\"{u}r die Functionen
$\displaystyle
   \frac{\partial u}{\partial x}
 - \frac{\partial \beta}{\partial y}$,
$\displaystyle
   \frac{\partial u}{\partial y}
 + \frac{\partial \beta}{\partial x}$
haben wir daher, wenn wir erstere $= X$ und letztere $= Y$
setzen, nicht bloss allgemein zu reden die Gleichung
\[ \frac{\partial X}{\partial x} + \frac{\partial Y}{\partial y}
   = 0,\]
sondern es wird auch durch die ganze Begrenzung irgend eines
Theils von $T$ erstreckt
\[ \int  \left(
            X \frac{\partial x}{\partial p}
          + Y \frac{\partial y}{\partial p}
         \right) \, ds = 0,\]
in so fern dieser Ausdruck \"{u}berhaupt einen bestimmten Werth
hat.

Zerlegen wir also (nach Art.~9, V) die Fl\"{a}che~$T$, wenn sie
einen mehrfachen Zusammenhang besitzt, durch Querschnitte in eine
einfach zusammenh\"{a}ngende $T^*$, so hat das Integral
\[ - \int_{O_0}^O
         \left(
            \frac{\partial u}{\partial p}
          + \frac{\partial \beta}{\partial s}
         \right) \, ds \]
f\"{u}r jede im Innern von $T^*$ von $O_0$ nach $O$ gehende Linie
denselben Werth und bildet, $O_0$ als fest gedacht, eine Function
von $x$,~$y$, welche in $T^*$ \"{u}berall eine stetige und
l\"{a}ngs eines Querschnitts beiderseits eine gleiche Aenderung
erleidet. Diese Function $\nu$ zu $\beta$ hinzugef\"{u}gt,
liefert uns eine Function $v = \beta + \nu$, von welcher der
Differentialquotient
$\displaystyle
      \frac{\partial v}{\partial x}
  = - \frac{\partial u}{\partial y}$
und
$\displaystyle
      \frac{\partial v}{\partial y}
  =   \frac{\partial u}{\partial x}$
ist.

Wir haben daher folgenden

\emph{Lehrsatz}.  Ist in einer zusammenh\"{a}ngenden, durch
Querschnitte in eine einfach zusammenh\"{a}ngende $T^*$ zerlegten
Fl\"{a}che $T$ eine complexe Function $\alpha + \beta i$ von
$x$,~$y$ gegeben, f\"{u}r welche
\[ \int
      \left[
         \left(
            \frac{\partial \alpha}{\partial x}
          - \frac{\partial \beta }{\partial y}
         \right)^2
       + \left(
            \frac{\partial \alpha}{\partial y}
          + \frac{\partial \beta }{\partial x}
         \right)^2
      \right] \, dT \]
durch die ganze Fl\"{a}che ausgedehnt einen endlichen Werth hat,
so kann sie immer und nur auf Eine Art in eine Function von $z$
verwandelt werden durch Hinzuf\"{u}gung einer Function
$\mu + \nu i$ von $x$,~$y$, welche folgenden Bedingungen
gen\"{u}gt:
\begin{description}

\item[\textmd{1)}]
$\mu$ ist am Rande $= 0$ oder doch nur in einzelnen Punkten davon
verschieden, $\nu$ in einem Punkte beliebig gegeben,

\item[\textmd{2)}]
die Aenderungen von $\mu$ sind in $T$, von $\nu$ in $T^*$ nur in
einzelnen Punkten und nur so unstetig, dass
\[ \int
      \left[
         \left( \frac{\partial \mu}{\partial x} \right)^2
       + \left( \frac{\partial \mu}{\partial y} \right)^2
      \right] \, dT
   \quad\mbox{und}\quad
   \int
      \left[
         \left( \frac{\partial \nu}{\partial x} \right)^2
       + \left( \frac{\partial \nu}{\partial y} \right)^2
      \right] \, dT \]
durch die ganze Fl\"{a}che erstreckt endlich bleiben, und
letztere l\"{a}ngs der Querschnitte beiderseits gleich.
\end{description}

Die Zul\"{a}nglichkeit der Bedingung zur Bestimmung von
$\mu + \nu i$ folgt daraus, dass $\mu$, durch welches $\nu$ bis
auf eine additive Constante bestimmt ist, stets zugleich ein
Minimum des Integrals $\Omega$ liefert, da, $u = \alpha + \mu$
gesetzt, offenbar f\"{u}r jedes $\lambda$ $N = 0$ wird; eine
Eigenschaft, die nach Art.~16 nur Einer Function zukommen kann.

\medbreak

\centerline{19.}

\nobreak\medskip

Die Principien, welche dem Lehrsatz am Schusse des vorigen
Art.\ zu Grunde liegen, er\"{o}ffnen den Weg, \emph{bestimmte}
Functionen einer ver\"{a}nderlichen complexen Gr\"{o}sse
(unabh\"{a}ngig von einem Ausdrucke f\"{u}r dieselben) zu
untersuchen.

Zu Orientirung auf diesem Felde wird ein Ueberschlag \"{u}ber den
Umfang der zur Bestimmung einer solchen Function innerhalb
eines gegebenen Gr\"{o}ssengebiets erforderlichen Bedingungen
dienen.

Halten wir uns zun\"{a}chst an einen bestimmten Fall, so kann,
wenn die \"{u}ber $A$ ausgebreitete Fl\"{a}che, durch welche dies
Gr\"{o}ssengebiet dargestellt wird, eine einfach
zusammenh\"{a}ngende ist, die Function $w = u + vi$ von $z$
folgenden Bedingungen gem\"{a}ss bestimmt werden:
\begin{description}

\item[\textmd{1)}]
f\"{u}r $u$ ist in allen Begrenzungspunkten ein Werth gegeben,
der sich f\"{u}r eine unendlich kleine Orts\"{a}nderung um eine
unendlich kleine Gr\"{o}sse von derselben Ordnung, \"{u}brigens
aber beliebig \"{a}ndert.\footnote{An sich sind die Aenderungen
dieses Werthes nur der Beschr\"{a}nkung unterworfen, nicht
l\"{a}ngs eines \emph{Theils} der Begrenzung unstetig zu sein;
eine weitere Beschr\"{a}nkung ist nur gemacht, um hier
unn\"{o}thige Weitl\"{a}ufigkeiten zu vermeiden.}

\item[\textmd{2)}]
der Werth von $v$ ist in irgend einem Punkte beliebig gegeben;

\item[\textmd{3)}]
die Function soll in allen Punkten endlich und stetig sein.
\end{description}

Durch diese Bedingungen aber ist sie vollkommen bestimmt.

In der That folgt dies aus dem Lehrsatze des vorigen Art., wenn
man, was immer m\"{o}glich sein wird, $\alpha + \beta i$ so
bestimmt, dass $\alpha$ am Rande dem gegebenen Werth gleich und
in der ganzen Fl\"{a}che f\"{u}r jede unendlich kleine
Orts\"{a}nderung die Aenderung von $\alpha + \beta i$ unendlich
klein von derselben Ordnung ist.

Es kann also, allgemein to reden, $u$ am Rande als eine ganz
willk\"{u}rliche Function von $s$ gegeben werden, und dadurch ist
$v$ \"{u}berall mit bestimmt; umgekehrt kann aber auch $v$ in
jedem Begrenzungspunkte beliebig angenommen werden, woraus dann
der Werth von $u$ folgt.  Der Spielraum f\"{u}r die Wahl der
Werthe von $w$ am Rande umfasst daher eine Mannigfaltigkeit von
Einer Dimension f\"{u}r jeden Begrenzungspunkt, und die
vollst\"{a}ndige Bestimmung derselben erfordert f\"{u}r jeden
Begrenzungspunkt Eine Gleichung, wobei es indess nicht wesentlich
sein wird, dass jede dieser Gleichungen sich auf den Werth Eines
Gliedes in Einem Begrenzungspunkte allein bezieht.  Es wird diese
Bestimmung auch so geschehen k\"{o}nnen, das f\"{u}r jeden
Begrenzungspunkt Eine mit der Lage dieses Punktes ihre Form
stetig \"{a}ndernde, beide Gleider enthaltende Gleichung gegeben
ist, oder f\"{u}r mehrere Theile der Begrenzung gleichzeitig so,
dass jedem Punkte eines dieser Theile $n - 1$ bestimmte Punkte,
aus jedem der \"{u}brigen Theile einer, zugesellt und f\"{u}r je
$n$ solcher Punkte gemeinschaftlich $n$ mit ihrer Lage stetig
ver\"{a}nderliche Gleichungen gegeben sind.  Diese Bedingungen,
deren Gesammtheit eine stetige Mannigfaltigkeit bildet und welche
durch Gleichungen zwischen willk\"{u}rlichen Functionen
ausgedr\"{u}ckt werden, werden aber, um f\"{u}r die Bestimmung
einer im Innern des Gr\"{o}ssengebiets \"{u}berall stetigen
Function zul\"{a}ssig und hinreichend zu sein, allgemein zu
reden, noch einer Beschr\"{a}nkung oder Erg\"{a}nzung durch
einzelne Bedingungsgleichungen---Gleichungen f\"{u}r
willk\"{u}rliche Constanten---bed\"{u}rfen, indem bis auf diese
sich die Genauigkeit unserer Sch\"{a}tzung offenbar nicht
erstreckt.

F\"{u}r den Fall, wo das Gebiet der Ver\"{a}nderlichkeit der
Gr\"{o}sse~$z$ durch eine mehrfach zusammenh\"{a}ngende
Fl\"{a}che dargestellt wird, erleiden diese Betrachtungen keine
wesentliche Ab\"{a}nderung, indem die Anwendung des Lehrsatzes in
Art.~18 eine bis auf die Aenderungen beim Ueberschreiten der
Querschnitte ebenso wie vorhin beschaffene Function
liefert---Aenderungen, welche $= 0$ gemacht werden k\"{o}nnen,
wenn die Grenzbedingungen eine der Anzahl der Querschnitte
gleiche Anzahl verf\"{u}gbarer Constanten enthalten.

Der Fall, wo im Innern l\"{a}ngs einer Linie auf Stetigkeit
verzichtet wird, ordnet sich dem vorigen unter, wenn man diese
Linie als einen Schnitt der Fl\"{a}che betrachtet.

Wenn endlich in einem einzelnen Punkte eine Verletzung der
Stetigkeit, also nach Art.~12 ein Unendlichwerden der Function,
zugelassen wird, so kann unter Beibehaltung der sonstigen in
unserm Anfangsfalle gemachten Voraussetzungen f\"{u}r diesen
Punkt eine Function von $z$, nach deren Subtraction die zu
bestimmende Function stetig werden soll, beliebig gegeben werden;
dadurch aber ist sie v\"{o}llig bestimmt.  Denn nimmt man die
Gr\"{o}sse $\alpha + \beta i$ in einem beliebig kleinen um den
Unstetigkeitspunkt beschriebenen Kreise gleich dieser gegebenen
Function, \"{u}brigens aber den fr\"{u}heren Vorschriften
gem\"{a}ss an, so wird das Integral
\[ \int
      \left(
         \left(
            \frac{\partial \alpha}{\partial x}
          - \frac{\partial \beta }{\partial y}
         \right)^2
       + \left(
            \frac{\partial \alpha}{\partial y}
          + \frac{\partial \beta }{\partial x}
         \right)^2
      \right) \, dT \]
\"{u}ber diesen Kreis erstreckt $= 0$, \"{u}ber den \"{u}brigen
Theil erstreckt einer endlichen Gr\"{o}sse gleich, und man kann
also den Lehrsatz des vorigen Art.\ anwenden, wodurch man eine
Function mit den verlangten Eigenschaften erh\"{a}lt.  Hieraus
kann man mit H\"{u}lfe des Lehrsatzes im Art.~13 folgern, dass im
Allgemeinen, wenn in einem einzelnen Unstetigkeitspunkte
die Function unendlich gross von der Ordnung $n$ werden darf,
eine Anzahl von $2n$ Constanten verf\"{u}gbar wird.

Geometrisch dargestellt liefert (nach Art.~15) eine Function~$w$
einer innerhalb eines gegebenen Gr\"{o}ssengebiets von zwei
Dimensionen ver\"{a}nderlichen complexen Gr\"{o}sse $z$ von einer
gegebenen $A$ bedeckenden Fl\"{a}che~$T$ ein ihr in den kleinsten
Theilen, einzelne Punkte ausgenommen, \"{a}hnliches, $B$
bedeckendes Abbild~$S$.  Die Bedingungen, welche so eben zur
Bestimmung der Function hinreichend und nothwendig befunden
worden sind, beziehen sich auf ihren Werth entweder in
Begrenzungs- oder in Unstetigkeitspunkten; sie erscheinen also
(Art.~15) s\"{a}mmtlich als Bedingungen f\"{u}r die Lage der
Begrenzung von $S$, und zwar geben sie f\"{u}r jeden
Begrenzungspunkt Eine Bedingungsgleichung.  Bezieht sich jede
derselben nur auf Einen Begrenzungspunkt, so werden sie durch
eine Schaar von Curven repr\"{a}sentirt, von denen f\"{u}r jeden
Begrenzungspunkt Eine den geometrischen Ort bildet.  Werden zwei
mit einander stetig fortr\"{u}ckende Begrenzungspunkte
gemeinschaftlich zwei Bedingungsgleichungen unterworfen, so
entsteht dadurch zwischen zwei Begrenzungstheilen eine solche
Abh\"{a}ngigkeit, dass, wenn die Lage des einen willk\"{u}rlich
angenommen wird, die Lage des andern daraus folgt.  Aehnlicher
Weise ergiebt sich f\"{u}r andere Formen der
Bedingungsgleichungen eine geometrische Bedeutung, was wir indess
nicht weiter verfolgen wollen.

\medbreak

\centerline{20.}

\nobreak\medskip

Die Einf\"{u}hrung der complexen Gr\"{o}ssen in die Mathematik
hat ihren Ursprung und n\"{a}chsten Zweck in der Theorie
einfacher\footnote{Wir betrachten hier als Elementaroperationen,
Addition und Subtraction, Multiplication und Division,
Integration und Differentiation, und ein Abh\"{a}ngigkeitsgesetz
als desto \emph{einfacher}, durch je weniger Elementaroperationen
die Abh\"{a}ngigkeit bedingt wird.  In der That lassen sich durch
eine endlich Anzahl dieser Operationen alle bis jetzt in der
Analysis benutzten Functionen definiren.}
durch Gr\"{o}ssenoperationen ausgedr\"{u}ckter
Abh\"{a}ngigkeitsgesetze zwischen ver\"{a}nderlichen Gr\"{o}ssen.
Wendet man n\"{a}mlich diese Abh\"{a}ngigkeitsgesetze in einem
erweiterten Umfange an, indem man den ver\"{a}nderlichen
Gr\"{o}ssen, auf welche sie sich beziehen, complexe Werthe giebt,
so tritt eine sonst versteckt bleibende Harmonie und
Regelm\"{a}ssigkeit hervor.  Die F\"{a}lle, in denen dies
geschehen ist, umfassen zwar bis jetzt erst ein kleines
Gebiet---sie lassen sich fast s\"{a}mmtlich auf diejenigen
Abh\"{a}ngigkeitsgesetze zwischen zwei ver\"{a}nderlichen
Gr\"{o}ssen zur\"{u}ckf\"{u}hren, wo die eine entweder eine
algebraische\footnote{D.~h.\ wo zwischen beiden eine algebraische
Gleichung Statt findet.}
Function der andern ist oder eine solche Function, deren
Differentialquotient eine algebraische Function ist---, aber
beinahe jede Schritt, der hier gethan ist, hat nicht bloss den
ohne H\"{u}lfe der complexen Gr\"{o}ssen gewonnenen Resultaten
eine einfachere, geschlossenere Gestalt gegeben, sondern auch zu
neuen Entdeckungen die Bahn gebrochen, wozu die Geschichte der
Untersuchungen \"{u}ber algebraische Functionen, Kreis- oder
Exponentialfunctionen, elliptische und Abel'sche Functionen den
Beleg liefert.

Es soll kurz andgedeutet werden, was durch unsere Untersuchung
f\"{u}r die Theories solcher Functionen gewonnen ist.

Die bisherigen Methoden, diese Functionen zu behandeln, legten
stets als Definition einen \emph{Ausdruck} der Function zu
Grunde, wodurch ihr Werth f\"{u}r \emph{jeden} Werth ihres
Argumentes gegeben wurde; durch unsere Untersuchung ist gezeigt,
dass, in Folge des allgemeinen Charakters einer Function einer
ver\"{a}nderlichen complexen Gr\"{o}sse, in einer Definition
dieser Art ein Theil der Bestimmungsst\"{u}cke eine Folge der
\"{u}brigen ist, und zwar ist der Umfang der
Bestimmungsst\"{u}cke auf die zur Bestimmung nothwendigen
zur\"{u}ckgef\"{u}hrt worden.  Dies vereinfacht die Behandlung
derselben wesentlich.  Um z.~B. die Gleichheit zweier
Ausdr\"{u}cke derselben Function zu beweisen, musste man sonst
den einen in den andern transformiren, d.~h.\ zeigen, dass beide
f\"{u}r jeden Werth der ver\"{a}nderlichen Gr\"{o}sse
\"{u}bereinstimmen; jetzt gen\"{u}gt der Nachweis ihrer
Uebereinstimmung in einem weit geringern Umfange.

Eine Theorie dieser Functionen auf den hier gelieferten
Grundlagen w\"{u}rde die Gestaltung der Function (d.~h.\ ihren
Werth f\"{u}r jeden Werth ihres Arguments) unabh\"{a}ngig von
einer Bestimmungsweise derselben durch Gr\"{o}ssenoperationen
festlegen, indem zu dem allgemeinen Begriffe einer Function einer
ver\"{a}nderlichen complexen Gr\"{o}sse nur die zur Bestimmung
der Function nothwendigen Merkmale hinzugef\"{u}gt w\"{u}rden,
und dann erst zu den verschiedenen Ausdr\"{u}cken deren die
Function f\"{a}hig ist \"{u}bergehen.  Der gemeinsame Charakter
einer Gattung von Functionen, welche auf \"{a}hnliche Art durch
Gr\"{o}ssenoperationen ausgedr\"{u}ckt werden, stellt sich dann
dar in der Form der ihnen auferlegten Grenz- und
Unstetigkeitsbedingungen.  Wird z.~B.\ das Gebiet der
Ver\"{a}nderlichkeit der Gr\"{o}sse~$z$ \"{u}ber die ganze
unendliche Ebene~$A$ einfach oder mehrfach erstreckt, und
innerhalb derselben der Function nur in einzelnen Punkten eine
Unstetigkeit, und zwar nur ein Unendlichwerden, dessen Ordnung
endlich ist, gestattet (wobei f\"{u}r ein unendliches $z$ diese
Gr\"{o}sse selbst, f\"{u}r jeden endlichen Werth $z'$ derselben
aber
$\displaystyle \frac{1}{z - z'}$
als ein unendlich Grosses erster Ordnung gilt), so ist die
Function nothwendig algebraisch, und umgekehrt erf\"{u}llt diese
Bedingung jede algebraische Function.

Die Ausf\"{u}hrung dieser Theorie, welche, wie bemerkt, einfache
durch Gr\"{o}ssenoperationen bedingte Abh\"{a}ngigkeitsgesetze
ins Licht zu setzen bestimmt ist, unterlassen wir indess jetzt,
da wir die Betrachtung des Ausdruckes einer Function
gegenw\"{a}rtig ausschliessen.

Aus demselben Grunde befassen wir uns hier auch nicht damit, die
Brauchbarkeit unserer S\"{a}tze als Grundlagen einer
\emph{allgemeinen} Theorie dieser Abh\"{a}ngigkeitsgesetze
darzuthun, wozu der Beweis erfordert wird, dass der hier zu
Grunde gelegte Begriff einer Function einer ver\"{a}nderlichen
complexen Gr\"{o}sse mit dem einer durch Gr\"{o}ssenoperationen
ausdr\"{u}ckbaren Abh\"{a}ngigkeit\footnote{Es wird darunter jede
durch eine endliche oder unendliche Anzahl der vier einfachsten
Rechnungsoperationen, Addition und Subtraction, Multplication und
Division, ausdr\"{u}ckbare Abh\"{a}ngigkeit begriffen.  Der
Ausdruck Gr\"{o}ssenoperationen soll (im Gegensatze zu
Zahlenoperationen) solche Rechnungsoperationen andeuten, bei
denen die Commensurabilit\"{a}t der Gr\"{o}ssen nicht in Betracht
kommt.}
v\"{o}llig zusammenf\"{a}llt.

\medbreak

\centerline{21.}

\nobreak\medskip

Es wird jedoch zur Erl\"{a}uterung unserer allgemeinen S\"{a}tze
ein ausgef\"{u}hrtes Beispiel ihrer Anwendung von Nutzen sein.

Die im vorigen Artikel bezeichnete Anwendung derselben ist,
obwohl die bei ihre Aufstellung zun\"{a}chst beabsichtigte, doch
nur eine specielle.  Denn wenn die Abh\"{a}ngigkeit durch eine
endliche Anzahl der dort als Elementaroperationen betrachteten
Gr\"{o}ssenoperationen bedingt ist, so enth\"{a}lt die Function
nur eine endliche Anzahl von Parametern, was f\"{u}r die Form
eines Systems von einander unabh\"{a}ngiger Grenz- und
Unstetigkeitsbedingungen, die zu ihrer Bestimmung hinreichen, der
Erfolg hat, dass unter ihnen l\"{a}ngs einer Linie in jedem
Punkte willk\"{u}rlich zu bestimmende Bedingungen gar nicht
vorkommen k\"{o}nnen.  F\"{u}r unsern jetzigen Zweck schien es
daher geeigneter, nicht ein dorther entnommenes Beispiel zu
w\"{a}hlen, sondern vielmehr ein solches, wo die Function der
complexen Ver\"{a}nderlichen von einer willk\"{u}rlichen Function
abh\"{a}ngt.

Zur Veranschaulichung und bequemeren Fassung geben wir demselben
die am Schlusse des Art.~19 gebrauchte geometrische Einkleidung.
Es erscheint dann als eine Untersuchung \"{u}ber die
M\"{o}glichkeit, von einer gegebenen Fl\"{a}che ein
zusammenh\"{a}ngendes in den kleinsten Theilen \"{a}hnliches
Abbild zu liefern, dessen Gestalt gegeben ist, wo also in obiger
Form ausgedr\"{u}ckt, f\"{u}r jeden Begrenzungspunkt des Abbildes
eine Ordscurve, und zwar f\"{u}r alle dieselbe, ausserdem aber
(Art.~5) der Sinn der Begrenzung und die Windungspunkte desselben
gegeben sind.  Wir beschr\"{a}nken uns auf die L\"{o}sung dieser
Aufgabe in dem Falle, wo jedem Punkte der einen Fl\"{a}che nur
Ein Punkt der andern entsprechen soll und die Fl\"{a}chen einfach
zusammenh\"{a}ngend sind, f\"{u}r welchen Fall sie in folgenden
Lehrsatz enthalten ist.

Zwei gegebene einfach zusammenh\"{a}ngende ebene Fl\"{a}chen
k\"{o}nnen stets so auf einander bezogen werden, dass jedem
Punkte der einen Ein mit ihm stetig fortr\"{u}ckender Punkt der
anderen entspricht und ihre entsprechenden kleinsten Theile
\"{a}hnlich sind; und zwar kann zu Einem innern Punkte und zu
Einem Begrenzungspunkte der entsprechende beliebig gegeben
werden; dadurch aber ist f\"{u}r all Punkte die Beziehung
bestimmt.

Wenn zwei Fl\"{a}chen $T$ und $R$ auf eine dritte $S$ so bezogen
sind, dass zwischen den entsprechenden kleinsten Theilen
Aehnlichkeit Statt findet, so ergiebt sich daraus eine Beziehung
zwischen den Fl\"{a}chen~$T$ und $R$, von welcher offenbar
dasselbe gilt.  Die Aufgabe, zwei beliebige Fl\"{a}chen auf
einander so zu beziehen, dass Aehnlichkeit in den kleinsten
Theilen Statt findet, ist dadurch auf die zur\"{u}ckgef\"{u}hrt,
jede beliebige Fl\"{a}che durch Eine bestimmte in den kleinsten
Theilen \"{a}hnlich abzubilden.  Wir haben hiernach, wenn wir in
der Ebene~$B$ um den Punkt, wo $w = 0$ ist, mit dem Radius~$1$
einen Kreis~$K$ beschreiben, um unsern Lehrsatz darzuthun, nur
n\"{o}thig zu beweisen: Eine beliebige einfach
zusammenh\"{a}ngende~$A$ bedeckende Fl\"{a}che~$T$ kann durch den
Kreis~$K$ stets zusammenh\"{a}ngend und in den kleinsten Theilen
\"{a}hnlich abgebildet werden und zwar nur auf Eine Art so, dass
dem Mittelpunkte ein beliebig gegebener innerer Punkt $O_0$ und
einem beliebig gegebenen Punkte der Peripherie ein beliebig
gegebeber Begrenzungspunkt $O'$ der Fl\"{a}che $T$ entspricht.

Wir bezeichnen die bestimmten Bedeutungen von $z$, $Q$ f\"{u}r
die Punkte $O_0$, $O'$ durch entsprechende Indices und
beschreiben in $T$ um $O_0$ als Mittelpunkt einen beliebigen
Kreis $\Theta$, welcher sich nicht bis zur Begrenzung von $T$
erstreckt und keinen Windungspunkt enth\"{a}lt.  F\"{u}hren wir
Polarcoordinaten ein, indem wir
$z - z_0 = r e^{\varphi i}$
setzen, so wird die Function
$\log (z - z_0) = \log r + \phi i$.
Der reelle Werth \"{a}ndert sich daher im ganzen Kreise mit
Ausnahme des Punktes $O_0$, wo er unendlich wird, stetig.  Der
imagin\"{a}re aber erh\"{a}lt, wenn \"{u}berall unter den
m\"{o}glichen Werthen von $\varphi$ der kleinste positive
gew\"{a}hlt wird, l\"{a}ngs des Radius, wo $z - z_0$ reelle
positive Werthe annimmt, auf der einen Seite den Werth~$0$, auf
der andern den Werth $2\pi$, \"{a}ndert sich aber dann in allen
\"{u}brigen Punkten stetig.  Offenbar kann dieser Radius durch
eine ganz beliebige vom Mittelpunkte nach der Peripherie gezogene
Linie~$l$ ersetzt werden, so dass die Function $\log (z - z_0)$
beim Uebertritt des Punktes $O$ von der negativen (d.~h.\ wo nach
Art.~8 $p$ negativ wird) auf die positive Seite dieser Linie eine
pl\"{o}tzliche Verminderung um $2\pi i$ erleidet, \"{u}brigens
aber sich mit dessen Lage im ganzen Kreise $\Theta$ stetig
\"{a}ndert.  Nehmen wir nun die complexe Function
$\alpha + \beta i$ von $x$,~$y$ im Kreise
$\Theta = \log (z - z_0)$, ausserhalb desselben aber, indem wir
$l$ beliebig bis an den Rand verl\"{a}ngern, so an, dass sie
\begin{description}
\item[\textmd{1)}]
an der Peripherie von $\Theta = \log (z - z_0)$, am Rande von $T$
bloss imagin\"{a}r wird,
\item[\textmd{2)}]
beim Uebertritt von der negativen auf die positive Seite der
Linie $l$ sich um $-2\pi i$, sonst aber bei jeder unendlich
kleinen Orts\"{a}nderung um eine unendlich kleine Gr\"{o}sse von
derselben Ordnung \"{a}ndert,
\end{description}
was immer m\"{o}glich sein wird, so erh\"{a}lt das Integral
\[ \int
      \left(
         \left(
            \frac{\partial \alpha}{\partial x}
          - \frac{\partial \beta }{\partial y}
         \right)^2
       + \left(
            \frac{\partial \alpha}{\partial y}
          + \frac{\partial \beta }{\partial x}
         \right)^2
      \right) \, dT,\]
\"{u}ber $\Theta$ ausgedehnt den Werth Null, \"{u}ber den ganzen
\"{u}brigen Theil erstreckt einen endlichen Werth, und es kann
daher $\alpha + \beta i$ durch Hinzuf\"{u}gung einer bis auf
einen bloss imagin\"{a}ren constanten Rest bestimmten stetigen
Function von $x$,~$y$, welche am Rande bloss imagin\"{a}r ist, in
eine Function $t = m + ni$ von $z$ verwandelt werden.  Der reelle
Theil $m$ dieser Function wird am Rande $= 0$, im Punkte
$O_0 = - \infty$ und \"{a}ndert sich im ganzen \"{u}brigen $T$
stetig.  F\"{u}r jeden zwischen $O$ und $-\infty$ liegenden Werth
$a$ von $m$ zerf\"{a}llt daher $T$ durch eine Linie, wo $m = a$
ist, in Theile, wo $m < a$ ist und die $O_0$ im Innern enthalten,
einerseits und anderseits in Theile, wo $m > a$ ist und deren
Begrenzung theils durch den Rand von $T$, theils durch Linien, wo
$m = a$ ist, gebildet wird.  Die Ordnung des Zusammenhangs der
Fl\"{a}che~$T$ wird durch diese Zerf\"{a}llung entweder nicht
ge\"{a}ndert oder erniedrigt, die Fl\"{a}che zerf\"{a}llt daher,
da diese Ordnung $= -1$ ist entweder in zwei St\"{u}cke von der
Ordnung des Zusammenhangs $0$ und $-1$, oder in mehr als zwei
St\"{u}cke.  Lezteres aber ist unm\"{o}glich, weil dann wenigstens
in Einem dieser St\"{u}cke $m$ \"{u}berall endlich und stetig und
in allen Theilen der Begrenzung constant sein m\"{u}sste,
folglich entweder in einem Fl\"{a}chentheil einen constanten
Werth, oder irgendwo---in einem Punkte oder l\"{a}ngs einer
Linie---einen Maximum- oder Minimumwerth haben m\"{u}sste, gegen
Art.~11, III.  Die Punkte, wo $m$ constant ist, bilden also in
sich zur\"{u}cklaufende allenthalben einfache Linien, welche ein
den Punkt $O_0$ einschliessendes St\"{u}ck begrenzen, und zwar
nimmt $m$ nach Innen zu nothwendig ab, woraus folgt, dass bei
einem positiven Umlaufe (wo nach Art.~8 $s$ w\"{a}chst) $n$
soweit es stetig ist, stets zunimmt, und also, da es nur beim
Uebertritt von der negativen auf die positive Seite der Linie~$l$
eine pl\"{o}tzliche Aenderung um $-2\pi$ erleidet\footnote{Da die
Linie~$l$ von einem im Innern des St\"{u}cks gelegenen Punkte bis
zu einem \"{a}ussern f\"{u}hrt, so muss sie, wenn sie dessen
Begrenzung mehrmals schneidet, Einmal mehr von Innen nach Aussen,
als von Aussen nach Innen gehen, und die Summe der
pl\"{o}tzlichen Aenderungen von $n$ w\"{a}hrend eines positiven
Umlaufs ist daher stets $= - 2 \pi$.}, jedem Werth zwischen $0$
und $2\pi$ Einmal von einem Vielfachen von $2\pi$ abgesehen
gleich wird.  Setzen wir nun $e^t = w$, so werden $e^m$ und $n$
Polarcoordinaten des Punktes~$Q$ in Bezug auf den Mittelpunkt des
Kreises~$K$.  Die Gesammtheit der Punkte~$Q$ bildet dann offenbar
eine \"{u}ber $K$ allenthalben einfach ausgebreitete
Fl\"{a}che~$S$; der Punkt $Q_0$ derselben f\"{a}llt auf den
Mittelpunkt des Kreises; der Punkt $Q'$ aber kann vermittelst der
in $n$ noch verf\"{u}gbaren Constante auf einen beliebig
gegebenen Punkt der Peripherie ger\"{u}ckt werden, w.~z.~b.~w.

In dem Falle, wo der Punkt $O_0$ ein Windungspunkt $(n-1) \,$\-ter
Ordnung ist, gelangt man, wenn nur $\log (z - z_0)$ durch
$\displaystyle \frac{1}{n} \log (z - z_0)$ ersetzt wird, durch
ganz \"{a}hnliche Schl\"{u}sse zum Ziele, deren weitere
Ausf\"{u}hrung man indess aus Art.~14 leicht erg\"{a}nzen wird.

\medbreak

\centerline{21.}

\nobreak\medskip

Die vollst\"{a}ndige Durchf\"{u}hrung der Untersuchung des
vorigen Artikels f\"{u}r den allgemeinern Fall, wo Einem Punkte
der einen Fl\"{a}che mehrere Punkte der andern entsprechen
sollen, und ein einfacher Zusammenhang f\"{u}r dieselben nicht
vorausgesetzt wird, unterlassen wir hier, zumal da, aus
geometrischem Gesichtspunkte aufgefasst, unsere ganze
Untersuchung sich in einer allgemeinern Gestalt h\"{a}tte
f\"{u}hren lassen.  Die Beschr\"{a}nkung auf ebene, einzelne
Punkte ausgenommen, schlichte Fl\"{a}chen, ist n\"{a}mlich
f\"{u}r dieselbe nicht wesentlich; vielmehr gestattet die
Aufgabe, eine beliebig gegebene Fl\"{a}che auf einer andern
beliebig gegebenen in den kleinsten Theilen \"{a}hnlich
abzubilden, eine ganz \"{a}hnliche Behandlung.  Wir begn\"{u}gen
uns, hier\"{u}ber auf zwei \emph{Gauss}'sche Abhandlungen, die
zu Art.~3 citirte und die disquis.\ gen.\ circa superf.\ art.~13,
zu verweisen.

\vfill\eject

\centerline{\large\textit{Inhalt}\footnote{[Diese
Inhalts\"{u}bersicht r\"{u}hrt fast vollst\"{a}ndig von Riemann her.]}.}

\nobreak\medskip

\begin{enumerate}
\item
Eine ver\"{a}nderliche complexe Gr\"{o}sse $w = u + vi$ heisst
eine Function einer andern ver\"{a}nderlichen Gr\"{o}sse
$z = x + yi$, wenn sie mit ihr sich so \"{a}ndert, dass
$\displaystyle \frac{dw}{dz}$
von $dz$ unabh\"{a}ngig ist.  Diese Definition wird begr\"{u}ndet
durch die Bemerkung, dass dies immer stattfindet, wenn die
Abh\"{a}ngigkeit der Gr\"{o}sse $w$ von $z$ durch einen
analytischen Ausdruck gegeben ist.

\item
Die Werthe der ver\"{a}nderlichen complexen Gr\"{o}ssne $z$ und
$w$ werden dargestellt durch die Punkte $O$ und $Q$ zweier Ebenen
$A$ und $B$, ihre Abh\"{a}ngigkeit von einander als eine
Abbildung der einen Ebene auf die andere.

\item
Ist die Abh\"{a}ngigkeit eine solche (Art.~1), dass
$\displaystyle \frac{dw}{dz}$
von $dz$ unabh\"{a}ngig ist, so findet zwischen dem Original und
seinem Bilde Aehnlichkeit in den kleinsten Theilen statt.

\item
Die Bedingung, dass
$\displaystyle \frac{dw}{dz}$
von $dz$ unabh\"{a}ngig ist, ist identisch mit folgenden:
$\displaystyle
   \frac{\partial u}{\partial x} =   \frac{\partial v}{\partial y}$,
$\displaystyle
   \frac{\partial u}{\partial y} = - \frac{\partial v}{\partial x}$.
Aus ihnen folgen
$\displaystyle
   \frac{\partial^2 u}{\partial x^2}
      + \frac{\partial^2 u}{\partial y^2}
   = 0$,
$\displaystyle
   \frac{\partial^2 v}{\partial x^2}
      + \frac{\partial^2 v}{\partial y^2}
   = 0$.

\item
Als Ort des Punktes $O$ wird f\"{u}r die Ebene $A$ eine begrenzte
\"{u}ber dieselbe ausgebreitete Fl\"{a}che~$T$ substituirt.
Windungspunkte dieser Fl\"{a}che.

\item
Ueber den Zusammenhang einer Fl\"{a}che.

\item
Das Integral
$\displaystyle \int
      \left(
         \frac{\partial X}{\partial x}
       + \frac{\partial Y}{\partial y}
      \right) \, dT$,
durch die ganze Fl\"{a}che $T$ erstreckt, ist gleich
$\displaystyle - \int
      (X \cos \xi + Y \cos \eta) \, ds$
durch ihre ganze Begrenzung, wenn $X$ und $Y$ beliebige in allen
Punkten von $T$ stetige Functionen von $x$ und $y$ sind.

\item
Einf\"{u}hrung der Coordinaten $s$ und $p$ des Punktes $O$ in
Bezug auf eine beliebige Linie.  Die gegenseitige
Abh\"{a}ngigkeit des Vorzeichens von $ds$ und $dp$ wird so
festgesetzt, dass
$\displaystyle
   \frac{\partial x}{\partial s} = \frac{\partial y}{\partial p}$
ist.

\item
Anwendung des Satzes im Art.~7, wenn in allen Fl\"{a}chentheilen
\[ \frac{\partial X}{\partial x} + \frac{\partial Y}{\partial y}
   = 0 \]
ist.

\item
Bedingungen, unter welchen im Innern einer $A$ einfach
bedeckenden Fl\"{a}che~$T$ eine Function~$u$, welche, allgemein
zu reden, der Gleichung
$\displaystyle
   \frac{\partial^2 u}{\partial x^2}
      + \frac{\partial^2 u}{\partial y^2}
   = 0$
gen\"{u}gt, nebst allen ihren Differentialquotienten \"{u}berall
endlich und stetig ist.

\item
Eigenschaften einer solchen Function.

\item
Bedingungen, unter welchen im Innern einer $A$ einfach
bedeckenden einfach zusammenh\"{a}ngenden Fl\"{a}che~$T$ eine
Function~$w$ von $z$ \"{u}berall nebst allen ihren
Differentialquotienten endlich und stetig ist.

\item
Unstetigkeiten einer solchen Function in einem inneren Punkte.

\item
Ausdehnung der S\"{a}tze der Art.~12 und 13 auf Punkte im Innern
einer beliebigen ebenen Fl\"{a}che.

\item
Allgemeine Eigenschaften der Abbildung einer in der Ebene~$A$
ausgebreiteten Fl\"{a}che~$T$ auf eine in der Ebene~$B$
ausgebreitete Fl\"{a}che~$S$, durch welche die Werthe einer
Function~$w$ von $z$ geometrisch dargestellt werden.

\item
Das Integral
$\displaystyle \int
      \left[
         \left(
            \frac{\partial \alpha}{\partial x}
          - \frac{\partial \beta }{\partial y}
         \right)^2
       + \left(
            \frac{\partial \alpha}{\partial y}
          + \frac{\partial \beta }{\partial x}
         \right)^2
      \right] \, dT$,
durch die ganze Fl\"{a}che~$T$ erstreckt, erh\"{a}lt bei
Aenderung von $\alpha$ um stetige oder doch nur in einzelnen
Punkten unstetige Functionen, die am Rande $= 0$ sind, immer
f\"{u}r Eine einem Minimumswerth und wenn man durch
Ab\"{a}nderung in einzelnen Punkten hebbare Unstetigkeiten
ausschliesst, nur f\"{u}r Eine.

\item
Begr\"{u}ndung eines in vorigen Art.\ vorausgesetzten Satzes
mittelst der Grenzmethode.

\item
Ist in einer beliebigen zusammenh\"{a}ngenden, durch Querschnitte
in eine einfach zusammenh\"{a}ngende $T^*$ zerlegten ebenen
Fl\"{a}che~$T$ eine Function $\alpha + \beta i$ von $x$,~$y$
gegeben, f\"{u}r welche
\[ \int
      \left[
         \left(
            \frac{\partial \alpha}{\partial x}
          - \frac{\partial \beta }{\partial y}
         \right)^2
       + \left(
            \frac{\partial \alpha}{\partial y}
          + \frac{\partial \beta }{\partial x}
         \right)^2
      \right] \, dT,\]
durch die ganze Fl\"{a}che erstreckt, endlich ist, so kann sie
immer und nur auf eine Art in eine Function von $z$ verwandelt
werden durch Hinzuf\"{u}gung einer Function $\mu + \nu i$ von
$x$,~$y$, welche so bedingt ist: 1) $\mu$ ist am Rande $= 0$,
$\nu$ in Einem Punkte gegeben.  2) Die Aenderungen von $\mu$ sind
in $T$, die von $\nu$ in $T^*$ nur in einzelnen Punkten und nur
so unstetig, dass
$\displaystyle \int
      \left[
         \left( \frac{\partial \mu}{\partial x} \right)^2
       + \left( \frac{\partial \mu}{\partial y} \right)^2
      \right] \, dT$
und
$\displaystyle \int
      \left[
         \left( \frac{\partial \nu}{\partial x} \right)^2
       + \left( \frac{\partial \nu}{\partial y} \right)^2
      \right] \, dT$
durch die ganze Fl\"{a}che endlich bleiben und letztere an den
Querschnitten beiderseits gleich.

\item
Ueberschlag \"{u}ber die hinreichenden und nothwendigen
Bedingungen zur Bestimmung einer Function complexen Arguments
innerhalb eines gegebenen Gr\"{o}ssengebiets.

\item
Die fr\"{u}here Bestimmungsweise eine Function durch
Gr\"{o}ssenoperationen enth\"{a}lt \"{u}berfl\"{u}ssige
Bestandtheile.  Durch die hier durchgef\"{u}hrten Betrachtungen
ist der Umfang der Bestimmungsst\"{u}cke einer Function auf das
nothwendige Mass zur\"{u}ckgef\"{u}hrt.

\item
Zwei gegebene einfach zusammenh\"{a}ngende Fl\"{a}chen k\"{o}nnen
stets so auf einander bezogen werden, dass jedem Punkte der einen
Ein mit ihm stetigfortr\"{u}ckender Punkt der andern entspricht
und ihre entsprechenden kleinsten Theile \"{a}hnlich sind; und
zwar kann zu Einem inneren Punkt und zu Einem Begrenzungspunkt
der entsprechende beliebig gegeben werden.  Dadurch ist f\"{u}r
alle Punkte die Beziehung bestimmt.

\item
Schlussbemerkungen.

\end{enumerate}

\end{document}

